\documentclass[11pt]{article}
\usepackage[margin=1.1in]{geometry}
\usepackage{amsmath,amssymb,amsthm,mathtools}
\usepackage{booktabs,graphicx,enumitem,url,xcolor}
\usepackage[colorlinks=true,linkcolor=blue!50!black,citecolor=blue!50!black,urlcolor=blue!50!black]{hyperref}
\newcommand{\E}{\mathbb{E}}
\newcommand{\PP}{\mathbb{P}}
\newcommand{\R}{\mathbb{R}}
\DeclareMathOperator{\Var}{Var}
\DeclareMathOperator{\Cov}{Cov}
\DeclareMathOperator{\Leb}{Leb}
\newcommand{\eqd}{\overset{d}{=}}
\newcommand{\one}{\mathbf{1}}
\newcommand{\Imm}{\mathcal{I}}
\newcommand{\Fc}{\mathcal{F}}
\newcommand{\Lc}{\mathcal{L}}
\newcommand{\hm}{\mathfrak{h}}
\newcommand{\Ac}{\mathcal{A}}
\newcommand{\Sc}{\mathsf{S}}
\theoremstyle{plain}
\newtheorem{theorem}{Theorem}[section]
\newtheorem{proposition}[theorem]{Proposition}
\newtheorem{lemma}[theorem]{Lemma}
\newtheorem{corollary}[theorem]{Corollary}
\newtheorem{claim}[theorem]{Claim}
\theoremstyle{definition}
\newtheorem{definition}[theorem]{Definition}
\newtheorem{ansatz}[theorem]{Ansatz}
\newtheorem{convention}[theorem]{Convention}
\newtheorem{constraint}{Constraint}
\newtheorem{principle}{Principle}
\newtheorem{example}[theorem]{Example}
\newtheorem{premise}{Premise}
\newtheorem{objection}[theorem]{Objection}
\renewcommand{\thepremise}{$\Phi$\arabic{premise}}
\theoremstyle{remark}
\newtheorem{remark}[theorem]{Remark}
\title{Technoanatomy:\\ \large Civilizational Survival as a Sign, Not a Size}
\author{Ajax Benander\thanks{Email:
\href{mailto:akb9408@rit.edu}{akb9408@rit.edu}. This paper develops the
existential-risk consequences of a survivorship model proved in a companion in
applied probability \cite{SurvA}, with the target-conditioning
machinery owned by the theorem companion \cite{TargetedB1}; both are
cited and restated, neither presupposed, and no downstream paper is
cited. Drafting
collaboration: Anthropic Claude (Opus~4.x, Fable~5), under the author's
direction.}}
\date{September 2026 --- Working Draft, rev.~4}
\begin{document}
\maketitle

\begin{abstract}
Existential-risk practice manages the wrong quantity. Its ledgers price
events, its progress measures index accumulated capacity, and both are
silent on the variable a survivorship model proves decisive: the sign of
what added capacity does to a civilization's internal hazard, here called
technoanatomy. The claim in one sentence: capability magnitude cannot substitute for
favorable anatomy, viability requiring favorable anatomy together with
favorable driver dynamics and a prerequisite schedule that is actually
cleared; on the schedule model entered here, whose leak law is
stipulated, the sign is being settled now at the highest leverage in the
record and conditioning on survival prefers that it be settled early;
and no event of bounded duration can certify which way it has gone. Four consequences
are drawn. The metric critique: an index of capability ranks worlds wrong
in both directions, and the Kardashev scale measures visibility, not
viability. The schedule: a third hazard channel, the gated sequence of
prerequisites a civilization must land, is formalized as a discounted
branching process feeding the survivorship model through a
surviving-completion kernel, with its derailment hazard, its modern-era
term acting on the anatomy rather than the capacity, a survival-time
threshold past which conditioning switches from favoring the unpushed
schedule to favoring the pushed one, survival-weighted toward better
completion states, and a
survivorship-prefers-early-change result under stated monotone-coupling
conditions; the leak law that makes the modern era the schedule's
high-leverage stage is stipulated, and every consequence that depends on
it is carried at that grade.
The rival periodization: the Precipice is engaged at its strongest and
three deltas are stated, in the quantity, the shape of time, and the
endpoint. And the two errors about the human role, extinctionism and
techno-optimism, refuted as mirror misreadings of one theorem. The
ceiling's institutional reading closes the paper: a governance built for
the regime where all the survivorship structure lives, with humility in
the objective rather than in the disclaimer, and no certified safety
anywhere.
\end{abstract}

\section{Introduction}\label{sec:intro}

The existential-risk field has, in a generation, built a serious
literature on how a civilization might end. It has not built an account
of what decides whether one can last, and the gap is not one of effort
but of quantity. The field's instruments price events: the probability
of a pandemic, a war, an unaligned system, summed and discounted into a
risk this century. Its progress measures index accumulated capacity:
energy captured, computation available, the Kardashev rungs. This paper
argues, from a theorem proved in a companion, that both instruments are
silent on the variable that decides survival, and that the silence is
structural rather than a matter of calibration. A civilization's added
capacity does two things at once: it muffles the hazards arriving from
outside, and it feeds a channel of hazard that arises from inside.
Whether indefinite survival is possible turns on the sign of the second
effect, not on the size of the first, and an index of capacity cannot
see the sign because the sign is not a function of what the index
measures.

\emph{The claim, in one sentence.} Capability magnitude cannot
substitute for favorable anatomy: viability requires a positive response
sign together with favorable driver dynamics and a prerequisite
schedule that is cleared, none of which a capacity index measures; on
the schedule model entered here, whose leak law is stipulated, the sign
is being settled now at the highest leverage in the record, conditioning
on survival prefers that it be settled early, and past a survival-time
threshold the same conditioning prefers the schedule pushed, tilted
toward better completion states; and no event of bounded duration can certify which way it
has gone.

\emph{What is imported.} The companion in applied probability
\cite{SurvA} proves the two-condition theorem on which the claim rests,
the model it is stated in, and the ceiling theorem that no finite
observation of a surviving path distinguishes the survivorship-conditioned
law from the unconditioned one. The theorem companion \cite{TargetedB1} supplies the general
machinery for records drawn under target-conditioned laws, used here
only where a result consumes it directly. The target this paper works with is the survivorship companion's own
viability event, the state past which the external channel is driven to
zero; no Golden-specific object is imported. A downstream companion on
the Route World hypothesis, not cited here so that the dependency runs
one way, gives that target its natural-historical reading and asks
whether it explains our record; this paper discusses that reading only
in its closing section. Each import is restated in
Section~\ref{sec:imports} at its grade; nothing later depends on more
than is stated there.
\emph{Position in the branch.} The Golden branch is four papers in a
strict order of import, and the citation graph runs the same way as the
logic. The survivorship companion \cite{SurvA} is the root and owns the
hazard model, the viability theorem, the half-time theorem, and the
hard-wall and soft-wall conditioned laws; the theorem companion
\cite{TargetedB1} owns the general result on records drawn under
target-conditioned laws and the model-comparison machinery, and imports
only the survivorship companion; this paper imports both and owns the
strength/anatomy distinction, the prerequisite schedule, and the
mechanism and policy consequences; and a terminal Route World companion
imports all three to state the hypothesis that the survivorship target
explains our record, citing this paper as a possible mechanism without
this paper citing it. Nothing here is offered as support for that
hypothesis, and nothing here assumes it.

\emph{What is new.} The reading of the viability theorem as a statement
about two state variables, strength and anatomy
(Section~\ref{sec:coords}). The schedule channel
(Section~\ref{sec:schedule}): the sequence of gated prerequisites a
civilization must land, formalized as a discounted branching process
with log-spread rates, its derailment hazard, its modern-era term acting
on the anatomy rather than on the capacity, the surviving-completion
kernel through which it feeds the survivorship model, the survival-time
threshold past which conditioning switches from favoring the unpushed
schedule to favoring the pushed one, tilted toward better completion states, and the front-loading
result under stated monotone-coupling conditions. The two applications
(Sections~\ref{sec:refutes} and \ref{sec:precipice}): the Kardashev
scale as the cleanest case of strength against anatomy, and the
Precipice as the cleanest comparison with existential-risk practice,
with three deltas stated. The remaining consequences are collected as a
discussion (Section~\ref{sec:discussion}).

\emph{What is refused.} No certified safety, anywhere, including the
paper's own favorable conclusions; no claim that the sign has been
settled favorably; no occupant of the decisive moment, which the
prefers-early result shows is felt at every moment. These are stated once
in Section~\ref{sec:fences}.

\emph{Grades.} [P] proved here or in a companion; [S] scaling-level or
structural, argued without proof; [$\Phi$] a stipulative commitment; [O]
open, with its specification.

\section{Vocabulary}\label{sec:vocab}
\begin{itemize}[itemsep=2pt]
\item \textbf{Technostrength} $x$: accumulated capacity, the integral of
the log-derivative driver; the quantity progress indices measure.
\item \textbf{Technoanatomy} $\rho$: the sign of what added capacity does
to the internal hazard; the decisive variable.
\item \textbf{Escape factor}: the driver's probability of escaping to
infinity; the second, sign-blind necessary condition.
\item \textbf{Departure / Golden Point}: the survivorship companion's
killing event and its complement, the state past which the external
hazard is driven to zero; labels here, not imports from downstream.
\item \textbf{Schedule}: the gated sequence of prerequisites, each a race
among competing clocks, that a civilization must land.
\item \textbf{Gate, clock, leak}: one prerequisite; one candidate route
through it; the decay of capacity that incompleteness imposes.
\item \textbf{Derailment hazard} $\sigma_t$: the rate at which a wrong
clock fires first at the open gate; the schedule's contribution to the
hazard of Departure.
\item \textbf{Anatomy-directed gate}: a gate whose resolution acts on
$\rho$ rather than on $x$; the modern era's kind.
\item \textbf{Settlement}: the point past which the survivorship-conditioned
mass of $\rho$ has largely moved; not a freezing.
\item \textbf{Leverage}: the sensitivity of the hazard to $\rho$, which
scales with $x$; why the anatomy was invisible in the deep past.
\item \textbf{The ceiling}: the imported theorem that no finite
observation certifies which measure a surviving record is under.
\item \textbf{Live third term}: the hypothesis, entered on evidence, that
the schedule is not yet exhausted.
\end{itemize}

\section{Imported results}\label{sec:imports}

\begin{definition}[The model; from \cite{SurvA}]\label{imp:model}
The total hazard is $h_t=B_t\varsigma_b(x_t)+\Gamma e^{-\rho x_t}$: an
external background $B$ muffled by a bounded logistic factor of the
accumulated capacity $x_t=\int_0^t\mu$, plus an internal term the same
capacity feeds or starves according to the sign of the response exponent
$\rho$; the driver $\mu$ is Brownian, so the noise sits on
$-\frac{d}{dt}\log h$. Departure is the killing time $T$; viability is
the event $\{\Lambda_\infty<\infty\}$ that the total hazard ever accrued
stays finite.
\end{definition}

\begin{definition}[The viability target; from \cite{SurvA}]\label{imp:goldenpoint}
The target is the survivorship companion's viability event, the state
past which the external channel is driven to zero and the total hazard
ever accrued stays finite; its complement is Departure. The names
``Golden Point'' and ``Departure Point'' used below are labels for these
two objects and carry no content beyond them here; the construction of
the Departure Point from the record's side, and the hypothesis that our
record is drawn under conditioning on this target, belong to the
downstream Route World companion and are not assumed.
\end{definition}

\begin{theorem}[Two conditions for viability; \textbf{[P]} in \cite{SurvA}]
\label{imp:twocond}
For $\rho\le0$ viability is null for every parameter value and every
driver. For $\rho>0$,
$\PP(\Lambda_\infty<\infty)=\one\{\rho>0\}\,V_{\mathrm{drv}}(\mu_0)$, with
$V_{\mathrm{drv}}$ the driver-viability functional, stated piecewise:
the scale-function escape probability for a transient driver, the
indicator of a positive stationary mean for a positive recurrent one,
zero for the Brownian driver, and open for general null recurrence. The
two factors are logically distinct necessary conditions, not independent
events: the first is a sign and the second is silent on it. No index of
accumulated capacity distinguishes the ascent that arrives from the
ascent that is a fall.
\end{theorem}

\begin{theorem}[Adverse credit occupies half of logarithmic time in the
critical regime; \textbf{[P]} in \cite{SurvA}]\label{imp:half}
For the driftless driver, the set of times at which the credit sits below
its initial value has logarithmic density exactly one half, almost
surely, for every initial datum; the statement is confined to the
critical recurrent regime and fails for any driver whose credit is
transient.
\end{theorem}

\begin{theorem}[The ceiling; \textbf{[P]} in \cite{SurvA}]\label{imp:ceiling}
In the critical regime the survivorship-conditioned law converges, in the
hard-wall model, to the Doob transform by the persistence-harmonic
function, $Q^{\mathrm{hard}}$, proved, with the soft-wall law
$Q^{\mathrm{soft}}$ conditional on the companion's ratio hypothesis;
either is locally absolutely continuous with respect to the
unconditioned law on survivors and globally singular. No finite observation of a surviving path certifies
which measure a record is under.
\end{theorem}

\section{Technostrength and technoanatomy}\label{sec:coords}

\subsection{Technostrength and technoanatomy}\label{gr:sub:coords}

\begin{definition}[The two coordinates]\label{gr:def:coords}
\emph{Technostrength} is the muffling exponent $x_t$: the capacity the
biosphere has accumulated, positive when its own organization holds its
hazard below background and negative when it holds it above.
\emph{Technoanatomy} is the exponent $\rho$ of
the survivorship companion \cite{SurvA}, the sign of
$-\partial_x\log h^{\mathrm{int}}$: whether acquiring capacity lowers or
raises the biosphere's hazard \emph{to itself}.

The two are of different kinds, and keeping them apart is the whole of
what follows. Technostrength is a \emph{state coordinate}: it
fluctuates, and by Theorem~\ref{imp:half} a biosphere in the
critical driftless regime, however well constituted, spends exactly half
of logarithmic time with its credit below its initial level. Technoanatomy is a \emph{structural property}: it does not
fluctuate at all, it says what capacity does when acquired, and
$\rho<0$ names the biosphere whose own strength is turned against it,
which is what the survivorship companion \cite{SurvA} means by eating oneself.
\emph{Adverse strength} is $x<0$; \emph{adverse anatomy} is $\rho<0$,
and only the second is a defect.
\end{definition}

Two results of Section~\ref{sec:imports} are worth restating in this vocabulary because
they are counterintuitive. \emph{Technostrength is not cumulative --- in the critical regime.}
Under a recurrent driver $x$ oscillates on scale $t^{3/2}$ and returns
to every level infinitely often: it is not a stock that accumulates but
a magnitude that fluctuates, so a framework treating it as monotonically
increasing has the wrong kind of object. Under a \emph{transient} driver
it does accumulate, and the ladder picture is locally right --- which is
the concession the broader Golden program makes precise and prices.

\emph{Every biosphere is adversely configured half the time.} By
Theorem~\ref{imp:half} the set $\{t:x_t<0\}$ has logarithmic density
exactly $\tfrac12$, for every parameter value and every biosphere.
Adverse \emph{strength} is therefore not a defect distinguishing some
biospheres from others but a condition all of them occupy half of
logarithmic time --- which is precisely why it cannot be what
``bad anatomy'' means, and why the two must be kept apart.

\subsection{What conditioning does to each coordinate}
\label{gr:subsub:measured}

\paragraph{Numerical methods.} All figures reported in this subsection,
in \S\ref{gr:subsub:anatsel}, in the broader Golden program, and in
the numerical check under the survivorship companion \cite{SurvA} were
produced by direct Monte Carlo on the two-channel model, using the exact
Gaussian update of the survivorship companion \cite{SurvA} for the driver-credit pair
and forward Euler for the background exponentials, at step
$\Delta=10^{-3}$. Conditioning on $\{T>t\}$ was performed by rejection, particles were killed by the Cox rule and the survivors reweighted
uniformly, so the conditioned estimates are $\PP^{(T)}$ estimates and
not $Q$ estimates, exactly as the column headings below say. Standard errors are binomial and are governed by the number of
\emph{survivors}, not of particles: at $T=4$ the survival probability
$2.2\times10^{-4}$ leaves about $26$ survivors of $1.2\times10^{5}$, so
the conditional probability $0.839$ carries a standard error near
$0.07$, and at $T=6$ about $23$ survivors give a comparable error; the
unconditioned entries, with $10^5$ effective samples, carry errors
below $0.004$. The conditional-anatomy figures are therefore read for
direction only, and a rerun by importance sampling, sequential Monte
Carlo, or splitting, with thousands of effective survivors, is owed
before any calibration is claimed \textbf{[O]}. \emph{Every $\PP(\rho>0\mid\text{survived})$
figure is prior-dependent}: $\rho\sim\mathrm{Unif}(-1,3)$ was chosen to
place a quarter of the mass on adverse anatomy and has no other
justification, and the $0.75\to0.84$ movement should be read as the
direction and rough magnitude of selection under that prior rather than
as a calibrated number. Code and seeds are available on request.

Two experiments are reported, and they are distinct. The strength rows
below are from the two-channel model of Section~\ref{sec:imports} at fixed
anatomy, $\sigma=1$, $\epsilon=0.05$, $n=3$, $b=3$, $\rho=2.5$, $k=6$,
conditioning horizon $T=4$, $N=1.2\times10^{5}$ particles. The anatomy
rows are from a separate run in which $\rho$ is drawn from the prior
$\mathrm{Unif}(-1,3)$ per particle and everything else is as above; its
survivor counts are the small ones just stated.

\begin{center}\small
\begin{tabular}{@{}lrr@{}}
\toprule
& General ($\PP$) & Conditioned ($\PP^{(4)}$)\\
\midrule
$\PP(x<0)$ (adverse strength) & $0.501$ & $0.0037$\\
median $x/\sigma T^{3/2}$ & $-0.001$ & $+1.085$\\
$\PP(\rho>0)$, $\rho\sim\mathrm{Unif}(-1,3)$ & $0.750$ & $0.839$\\
$5\%$ quantile of $\rho$ & $-0.80$ & $-0.44$\\
\bottomrule
\end{tabular}
\end{center}

\noindent The two $\rho$ rows depend on the prior $\mathrm{Unif}(-1,3)$
placed on anatomy, which is a choice and not a result; the $x$ rows do
not.

\begin{remark}[Truncation versus rescaling]\label{gr:rem:truncresc}
Two numbers, and they differ in kind rather than in force.
Adverse strength is suppressed by a factor of $137$ --- abruptly, and
within a few multiples of the intrinsic timescale. Technoanatomy is
selected too, but \emph{gradually}: $\PP(\rho>0)$ moves from $0.75$ to
$0.839$ at the same horizon, reaching $0.909$ only by $T=6$, with the
negative tail truncating steadily ($5\%$ quantile $-0.80\to-0.44\to
-0.14$).

The asymmetry is structural rather than incidental, and it follows from
what the two objects are. Strength is a state coordinate, so
conditioning pushes it out of a half-space it can and does re-enter ---
fast, and reversibly. Anatomy is a fixed structural parameter, so it can
be selected only through the accumulated hazard cost of never changing
--- slow, and irreversibly. Anatomy is the slower and the more decisive:
by Theorem~\ref{imp:twocond}, adverse anatomy is eventually fatal
with certainty, while adverse strength is a condition a biosphere in the critical
recurrent regime occupies half of logarithmic time and survives.

\end{remark}

\subsection{Does conditioning select technoanatomy?}
\label{gr:subsub:anatsel}

The measurement above concerns technostrength, a state coordinate; the
claim the argument needs is about technoanatomy, a structural parameter,
which therefore cannot be selected \emph{within} a trajectory. It can be
selected across an ensemble, and that is directly computable: place a
prior on $\rho$, evolve, and reweight by survival. Taking
$\rho\sim\mathrm{Uniform}(-1,3)$, so that $\PP(\rho>0)=0.75$ a priori:

\begin{center}\small
\begin{tabular}{@{}rrrr@{}}
\toprule
$T$ & $\PP(\text{survive})$ & $\PP(\rho>0\mid\text{survived})$ &
$5\%$ quantile of $\rho$\\
\midrule
$2$ & $2.8\times10^{-3}$ & $0.745$ & $-0.80$\\
$4$ & $2.2\times10^{-4}$ & $0.839$ & $-0.44$\\
$6$ & $1.9\times10^{-4}$ & $0.909$ & $-0.14$\\
\bottomrule
\end{tabular}
\end{center}

\noindent Survivorship does select for positive technoanatomy, and the
adverse tail truncates steadily. But the selection is \emph{gradual}
where the strength truncation was abrupt: adverse strength falls by a
factor of $137$ at $T=4$, while $\PP(\rho>0)$ has moved only from
$0.75$ to $0.84$ at the same horizon. Anatomy is the slower and the more
decisive: by Theorem~\ref{imp:twocond} an adverse $\rho$ is
eventually fatal with certainty, while adverse strength is a condition
a biosphere in the critical regime occupies half of logarithmic time and
survives.

\paragraph{The coordinates stay separable.}\label{gr:rem:separable}
Strength and anatomy cannot fuse under conditioning, and the reason is
structural rather than measured: $\rho$ is a parameter of the model and
$x$ a functional of the path, so no amount of conditioning makes one a
function of the other. They therefore remain distinct diagnostic
questions rather than two aspects of a single ``advancement'' --- which
is exactly what a one-dimensional scale cannot represent.



\section{The third channel: the schedule}\label{sec:schedule}


Definition~\ref{imp:model} carries two channels, and
the downstream Route World hypothesis records that they do not exhaust the
hazard of departure: $h^{\mathrm{tot}}$ prices $\mathrm{Base}$ and
nothing else, while the hazard of departure is the hazard of $\mu G$,
which by the downstream Route World hypothesis is at least as large and in general
strictly larger. That remark calls the larger rate open \textbf{[O]}. This subsection is a first attempt at it. The claim is modest: the rate
has a computable shape, the shape is a third channel, and the
distribution it needs is one the companion already supplies.

\subsection{A collision of vocabulary, settled first}
\label{sch:subsub:vocab}

\begin{convention}[Endogenous, internal, schedule]\label{sch:conv:internal}
The companion divides departure into \emph{external}, inflicted from
outside the biosphere, and \emph{internal}, defined there as the
complement: the remaining cases, all occurring within the biosphere's
scope \cite{Nexus}. That category is a residue and is correspondingly
wide; the companion's own instance of it is a biosphere too
underdeveloped to reach a milestone given its resources and timing,
which is a scheduling failure and not a self-inflicted one.

The internal channel of Definition~\ref{imp:model} is narrower: hazard the
biosphere inflicts on itself with capacity it has acquired, whose
response to that capacity is the technoanatomy $\rho$ of
the survivorship companion \cite{SurvA}.

The two senses are a category and one of its members. This work uses
\emph{endogenous} for the companion's residue and reserves
\emph{internal} for the channel:
\begin{equation}\label{sch:eq:partition}
\text{endogenous}\;=\;\underbrace{\text{internal}}_{\text{capacity
turned inward}}\;+\;\underbrace{\text{schedule}}_{\text{derailment
before capacity can be held}}\;+\;\cdots
\end{equation}
The ellipsis is substantive. A category defined by exclusion is not
shown exhausted by exhibiting two members, and nothing here claims these
are the last.\footnote{The rename is a convention of this work and is
offered rather than imposed; \emph{non-external} and \emph{endemic}
would serve. One point of order regardless of the choice: the companion
is deposited, so a reader arriving from the deposit will meet
\emph{internal doomsday} in the residue sense, and the convention is
owed to bridge the two states of that text whether or not the term is
changed there.}
\end{convention}

\subsection{The schedule}\label{sch:subsub:model}

\begin{definition}[Schedule; \textbf{[S]}]\label{sch:def:schedule}
A \emph{schedule} is a sequence of $n$ gates to be passed in order. At
gate $k$ there are $m$ competing exponential clocks with rates
$\lambda_{k,1},\dots,\lambda_{k,m}$, the first corresponding to the
realized option. Write
\begin{equation}\label{sch:eq:clocks}
\lambda_k=\sum_{j=1}^{m}\lambda_{k,j},
\qquad
\pi_k=\frac{\lambda_{k,1}}{\lambda_k},
\end{equation}
so that gate $k$ fires at rate $\lambda_k$ and resolves to the realized
option with probability $\pi_k$. Write $k_t$ for the number of gates
passed by time $t$. The \emph{onset state} is $k=n$: the first state at
which the muffling exponent $x$ of the survivorship companion \cite{SurvA} can be held at all.
\end{definition}

The rates are the primitives and $\lambda_k,\pi_k$ are read off them,
which matters because it is the per-option rates a natural-historical
account would estimate --- not a tempo and a fallibility chosen
separately. Stage dependence is retained deliberately. The reason the object is a
\emph{sequence} rather than a set is that each event's effect depends on
the state its predecessor left: the climatic and ecological
baton-passing that carried a lineage down from the trees is not a set of
independent draws, and a constant $\pi$ would model it as one.
Stage-dependent $\pi_k$ is the minimum that preserves the structure;
state-dependent $\pi(\cdot)$ is what the picture actually asks for, and
is not supplied here \textbf{[O]}.

One point of scope, since it is a natural worry and the construction
already absorbs it. An event later in the sequence arriving before its
predecessors needs no transition of its own: it is one of the $m$
options competing at the gate it arrives at, and resolving to it is a
wrong resolution like any other. The $m-1$ non-realized options are not
a homogeneous class, each carries its own rate by
Definition~\ref{sch:def:rates} and its own consequence, so
heterogeneity of failure is present without being labeled by type. What
the chain does not carry is partial progress within a gate: it is Markov
between them, which is a simplification and is chosen rather than
derived.

The discreteness is a limitation of the same kind as calling Earth a
sphere. Exotic structure in the ordering is not recoverable from a chain
over $m^n$ articulable realizations; the shape of the dependence is, and
that is what is wanted here.

\subsection{Where the rates come from}\label{sch:subsub:rates}

Definition~\ref{sch:def:schedule} takes the per-option rates as given.
One specification of how they are drawn is worth carrying, because it
determines more than it looks like it should.

\begin{definition}[Log-spread rates; \textbf{[S]}]\label{sch:def:rates}
At each gate the $m$ rates are drawn afresh and independently with
$Y_{k,j}:=\log\lambda_{k,j}$ i.i.d.\ from a \emph{symmetric log-spread
family} $F$: a law with finite variance and tails tempered as the
application requires, a shortest feasible time bounding the fast end
and a longest tolerable one the slow; and \emph{which option is the
realized one is independent of its rate}. The correct option may be the
fastest of the $m$, or the slowest, and nothing about the draw favors
either. Every consequence below uses only that the $Y_{k,j}$ are
i.i.d.\ under $F$, so that a log-rate gap $Y_i-Y_j$ is symmetric about
zero. The canonical instance is $F=N(\mu,\sigma^2)$, the lognormal
rate, which carries its own justification: a log-rate that is a sum of
many independent multiplicative contributions is normal by the central
limit theorem, so the law is inherited rather than chosen.
\end{definition}

Three consequences, in ascending order of how much they change. \emph{The waiting time is set by the correct option's own rate.} A gate
fires at $\lambda_k=\sum_j\lambda_{k,j}$ and resolves correctly with
probability $\pi_k=\lambda_{k,1}/\lambda_k$, so the expected number of
attempts before a correct resolution is $1/\pi_k$ and each costs
$1/\lambda_k$. The product is $1/\lambda_{k,1}$: however many wrong
options crowd the gate, what a biosphere waits on is the realized
option's own clock. Exposure per gate is therefore log-spread along with
the rate, and the difference between a lucky gate and an unlucky one is
orders of magnitude of waiting rather than a factor.

\emph{The contrast between options is a random order of magnitude, with
the sign withheld.} $\log(\lambda_{k,1}/\lambda_{k,j})=Y_{k,1}-Y_{k,j}$ is a
difference of two i.i.d.\ draws from $F$: symmetric about zero, spread
across the range, and committing to no direction. That is the shape of the High-Contrast
Earth Hypothesis, and the correspondence is exact rather than
suggestive --- the Hypothesis asserts a random order of magnitude between
independently constructed alternatives and \emph{deliberately withholds
the sign}, with the Strong Accordance Principle supplying direction
separately. \eqref{sch:def:rates} supplies the same shape from a
mechanical premise about competing clocks, and the direction arrives the
same way it does there: from conditioning, below. So the Hypothesis
acquires a microfoundation at the gate level, and does not have to be
imported as a posit at this scale --- though what replaces it is another
\textbf{[S]}, and the trade is one posit for a more checkable one rather
than a derivation from nothing.

The correspondence should be stated at the right strength, since it is
not an identity of distributions. By \eqref{sch:eq:cefff} the contrast a
perturbation carries is $r\log_{10}(1+(B+\Gamma)/\lambda)$, which in the
regime supplying many orders (Remark~\ref{sch:rem:tension}) is $r$ times
a \emph{log-rate gap} --- and under Definition~\ref{sch:def:rates} a
log-rate gap is a difference of i.i.d.\ draws from $F$, hence symmetric
and spread across orders with no committed sign. The companion states
the Hypothesis as uniform rather than normal. What the collapse
consumes, however, is not the shape but two functionals of it, the
typical contrast $\E[c]$ and the averaged one
$-\log_{10}\E[10^{-c}]$, and both are defined for either law. So the
construction is compatible with the Hypothesis in the sense the argument
uses it, without reproducing its stated distribution.

The scale of the spread is a further matter, and the accurate thing to say is not that it awaits measurement. Nothing available to an observer
inside a biosphere fixes what \emph{sets} the range over which competing
rates are distributed: one meets the distribution already running, with
no vantage from which its parameters could be explained rather than
merely estimated. That is Theorem~\ref{imp:ceiling}'s structure applied one
level down, and the scale is scoped as uncertifiable for the reason
$\rho$ is, rather than listed as a parameter someone has yet to
supply.

\emph{And conditioning selects the fast-correct branch, which is why the
realized path looks optimal.} Exposure at gate $k$ is $1/\lambda_{k,1}$,
paid at hazard $B+\Gamma$ by Proposition~\ref{sch:prop:gating}, so a
branch's survival weight carries $\exp\{-(B+\Gamma)/\lambda_{k,1}\}$ at
every gate. Branches in which the realized option happened to be slow
pay exponentially for it. Conditioning on survival therefore selects,
gate by gate, for histories in which the correct option was among the
fastest available --- not because anything chose them, but because the
alternatives were charged for the wait.

That is the mechanism, and should be stated as one. Under the Cosmic Nexus the
realized option is a random draw and is usually slow; under conditioning
one is overwhelmingly likely to find oneself on a branch where it was
fast at gate after gate. The appearance of an optimal path threaded
through a combinatorial mess is what survivorship does to a log-spread
rate landscape, and requires no optimization to produce it. This is the
Optimal-Earth phenomenology arriving from the schedule's own machinery
rather than being assumed alongside it \cite{Nexus, Ben26Optimal}.

\subsection{Incompleteness decays capacity}\label{sch:subsub:leak}

The first commitment is that until the schedule completes, the muffling
exponent cannot be held. This departs from
Definition~\ref{imp:model}, which states that there is no capacity
parameter and no leak; the departure is deliberate and is confined.

\begin{definition}[Staged leak; \textbf{[S]}]\label{sch:def:leak}
Replace the survivorship companion \cite{SurvA} and the survivorship companion \cite{SurvA} by
\begin{equation}\label{sch:eq:leak}
dx_t=\mu_t\,dt-\gamma_{k_t}x_t\,dt,
\qquad
d\mu_t=-\nu_{k_t}\mu_t\,dt+\sigma\,dW_t,
\end{equation}
with $\gamma_k,\nu_k$ nonincreasing in $k$ and $\gamma_n=\nu_n=0$. The
schedule of coefficients is \emph{geometric}: each closed gate
multiplies the leak by a constant factor,
\begin{equation}\label{sch:eq:geomleak}
\gamma_k=\gamma_0\,\varphi^{\,k},\qquad
\varphi=(\gamma_{n-1}/\gamma_0)^{1/(n-1)}<1,
\end{equation}
with $\gamma_0$ enormous, so that before any gate is closed the credit
is pinned at the origin to a degree with no practical upper bound, and $\gamma_n=0$ at completion.
\end{definition}

The solution of the first equation, with the stage index varying along
the path, is
\[
x_t=e^{-\int_0^t\gamma_{k_u}du}\,x_0+\int_0^te^{-\int_s^t\gamma_{k_u}du}\,\mu_s\,ds,
\]
which reduces to the convolution $\int_0^te^{-\gamma_k(t-s)}\mu_s\,ds$
only conditional on remaining at one stage $k$ throughout $[0,t]$: the
credit is still the integral of the driver, accumulated additively,
with each increment buffered by a multiplicative decay from the moment
it lands, the decay rate being whatever stage is open when it lands. The leak is
not an alternative to the survivorship companion \cite{SurvA} but a discounting of it. The second equation is the same commitment one derivative up, and it is
what the picture requires rather than an extra assumption. The driver
$\mu=-\tfrac{d}{dt}\log h$ is an \emph{improvement rate}, and a rate of
improvement is something an agent supplies. Before the onset state there
is none, so $\mu$ is not a free random walk either: neither the stock nor
the rate is sustained, and the two dampings are one absence expressed at
two levels. A biosphere's ancestors do not hold a small amount of
technostrength --- they hold none, and what the pair of equations says is
that excursions away from zero are transient rather than accumulated.

At stage $k$, with the forcing $\mu$ switched off, the homogeneous part
obeys $\tfrac{d}{dt}\log x=-\gamma_k$: capacity acquired before
completion bleeds away at a rate each completed gate reduces, the
forced process adding the driver's increments on top of that decay. The leak is by hand, and making it emergent, $\gamma$
falling out of the schedule's own structure rather than being assigned
to it, is the model's natural next target \textbf{[O]}.

\paragraph{What a closed gate buys, and why almost none of it is
bought early.}\label{sch:rem:backloaded}
\eqref{sch:eq:geomleak} makes gates equal in $\log\gamma$ rather than in
$\gamma$, and the consequence for the payoff is severe enough to state
on its own. The credit a biosphere can sustain at stage $k$ scales as
$1/\gamma_k=\varphi^{-k}/\gamma_0$, so it grows \emph{geometrically} in
gates closed. But the hazard depends on the credit \emph{exponentially}, both channels of Definition~\ref{imp:model} do, by construction, so the
relief a closed gate delivers goes as
$\exp\{-\rho C\varphi^{-k}/\gamma_0\}$, doubly exponential in $k$.
Closing the last gate is worth more than closing all the others
together, by a margin that is not close.

With $\gamma_0$ enormous the reading sharpens. The early gates take a
biosphere from a credit that is unimaginably pinned to one merely very
pinned, which is no relief at all in the only currency the hazard reads.
Practical relief arrives in the final stretch, and a schedule
half-climbed is nearer in hazard to one not begun than to one completed. Set this beside the broader Golden program, which has stakes rising
monotonically for combinatorial reasons --- a late step carries the
improbability of every step behind it. The present observation is a
second and independent reason for the same monotonicity, arriving from
the leak's geometry rather than the chain's combinatorics. Two arguments
of different provenance agreeing on the profile is a better position
than either alone.

\begin{remark}[The existing model is the completed case]
\label{sch:rem:limit}
$\gamma_n=0$ is not a convenience. It makes
Definition~\ref{imp:model} the $k=n$ boundary case of
\eqref{sch:eq:leak} exactly, so nothing proved in
Section~\ref{sec:imports} is disturbed: the
persistence exponent $\tfrac14$, the trichotomy, the two-condition
theorem and the survivorship interpolation are all statements about the
completed regime, and this channel prefixes them rather than amending
them.

The price is that the incomplete regime has its own asymptotics --- and
they can be given, which is the next proposition's business. The
persistence problem is not the one the survivorship companion \cite{SurvA} solves, and
its exponent does not carry.
\end{remark}

\begin{proposition}[The incomplete regime has no polynomial tail;
\textbf{[S]}]\label{sch:prop:exponent}
Under \eqref{sch:eq:leak} with $\gamma_k,\nu_k>0$ held fixed, the pair
$(\mu,x)$ is a two-dimensional linear system with both eigenvalues
negative, driven by white noise; it is therefore ergodic with a
stationary Gaussian law centered at the origin. Persistence of $x$ is then
persistence of a stationary Gaussian process, which decays
\emph{exponentially}:
\begin{equation}\label{sch:eq:incompletetail}
\PP\bigl(x\ \text{stays positive to }t\bigr)\;\asymp\;e^{-\theta(\gamma_k,\nu_k)\,t},
\qquad \theta>0 .
\end{equation}
No power law survives. The rate $\theta$ is not given in closed form
here; persistence rates for stationary Gaussian processes are generally
not, and \eqref{sch:eq:incompletetail} identifies the class rather than
the constant.\footnote{Simulation at $\sigma=1$, $8\times10^{4}$ paths.
With $\gamma=1$ and the driver \emph{undamped} ($\nu=0$) the decay is a
power law, fitted slope $-0.476$, log--log correlation $0.9995$, which is the damped randomly accelerated particle and is not the regime
this definition describes. With the driver damped, the fit inverts
decisively: at $(\gamma,\nu)=(1,0.5)$ and $(1,1)$ the correlations
against $t$ are $1.0000$ and $0.9999$ with rates $0.222$ and $0.312$,
against log--log correlations of $0.975$ and $0.985$.}
\end{proposition}

\paragraph{What incompleteness costs, and what completion buys.}\label{sch:rem:suspends}
The contrast with the completed regime is not a change of exponent but a
change of kind, and it is the sharpest thing the channel says about why
the schedule is worth climbing.

Section~\ref{sec:imports} establishes that the completed model departs with
probability one and yet has infinite life expectancy: the tail is
$t^{-1/4}$, residual life is Pareto, quantiles grow as the fourth power
of rarity, and the effective hazard is $1/4t$ --- falling at every age.
By \eqref{sch:eq:incompletetail} the incomplete regime has none of it.
An exponential tail has finite expectation, memoryless residual life,
quantiles growing logarithmically, and a hazard that does not fall at
all. \emph{Every consequence of the survivorship companion \cite{SurvA} is a
consequence of having completed}, and none of it is available before.

Two readings follow, and the second is structural. The Lindy structure is a \emph{reward} rather than a property. A
biosphere before its onset state is not on a fat tail at low altitude;
it is on a thin one, and no amount of waiting improves its position,
since an exponential is memoryless. What completing the schedule buys is
not more of the same but a different tail entirely.

And the floor is the trapped regime's floor with one difference that
matters. Theorem~\ref{imp:twocond}(ii) gives $\PP(T>t)\le e^{-h_\ast t}$
for adverse anatomy at frozen background, and \eqref{sch:eq:incompletetail}
gives an
exponential floor for incompleteness. The two are structurally alike and
disposed of oppositely: the anatomy floor cannot be left, since $\rho$ is
a parameter, while the schedule floor is \emph{closable} --- $\gamma_k$
and $\nu_k$ vanish at completion, and the debt that stands between is
payable whenever the criterion of \S\ref{sch:subsub:open} holds. An
incomplete biosphere and a trapped one sit on the same kind of tail; only
one of them can get off it.

A biosphere may of course be both, and the two floors then compose
exactly rather than competing. If $h^{\mathrm{tot}}\ge h_\ast$ then
$\PP(T>t)=e^{-h_\ast t}\,\E\bigl[e^{-\int(h^{\mathrm{tot}}-h_\ast)}\bigr]$
identically, a constant floor factors out of the exponential
functional, leaving the damped-persistence problem for the recentered
hazard with its own rate $\tilde\theta$. So
$\theta_{\mathrm{joint}}=h_\ast+\tilde\theta$: neither floor governs, they
\emph{add}, and the interventions are correspondingly separable.
Completing the schedule removes $\tilde\theta$ and not $h_\ast$; moving
$\rho$ removes $h_\ast$ and not $\tilde\theta$. A biosphere that is both
trapped and incomplete must do both things, and doing either buys exactly
its own term.

\subsection{The schedule gates the external channel}
\label{sch:subsub:gating}

Completing the chain confers no capacity. It confers only the ability to
\emph{hold} capacity --- which sounds like a small thing and is not,
because of what Definition~\ref{imp:model} does at $x=0$.

\begin{proposition}[An unheld $x$ leaves the crest at bare background;
\textbf{[P]} in the pinned schedule model]\label{sch:prop:gating}
Two models of the incomplete state are used in this section and are
kept apart. In the \emph{pinned schedule model} the credit is
identically zero while the schedule is incomplete, $x\equiv0$; in the
\emph{leaky schedule refinement} of Definition~\ref{sch:def:leak} the
incomplete credit is a damped stochastic process, $dx=\mu\,dt-\gamma_kx\,dt$.
This proposition, the completion-weight fixed point
\eqref{sch:eq:fixedpoint}, and the threshold proposition below are
exact for the pinned model and are surrogate, scaling-level statements
for the leaky refinement, in which the incomplete hazard is
$B_t+\Gamma e^{-\rho x_t}$ with $x_t$ small rather than zero. In the
pinned model, by the normalization of Definition~\ref{imp:model},
$\varsigma_b(0)=1$ exactly, so $h^{\mathrm{crest}}\big|_{x=0}=B_t$; and the trough does not
vanish there either, since $\Gamma e^{-\rho\cdot0}=\Gamma$. A biosphere
that cannot hold $x$ therefore faces the external background unshielded
\emph{and} pays the internal channel's full coefficient, at total
$h^{\mathrm{tot}}\big|_{x=0}=B_t+\Gamma$. Its cumulative hazard is
$\Lambda_t=\int_0^t(B_u+\Gamma)\,du\to\infty$, and
$\PP(\Lambda_\infty<\infty)=0$ identically --- for every parameter
value, every anatomy $\rho$, and every driver. The configuration and its
consequences are those of the broader Golden program.
\end{proposition}

\begin{remark}[The external term is small only because the schedule
succeeded]\label{sch:rem:reliance}
Read at $x>0$ the crest channel is the least of the model's worries: it
is bounded above by the survivorship companion \cite{SurvA}, it is shielded
without limit as capacity accumulates, and beside an unbounded internal
channel it looks like one measly term. That reading presupposes what the
schedule supplies. Run the dependency out. Shielding is the factor $\varsigma_b(x)$;
$\varsigma_b$ shields only for $x>0$; $x$ can be held only at $k=n$; and
$k=n$ is reached through $n$ gates each of which can go $m-1$ ways
wrong. \emph{So the survivability of the external channel is downstream
of the schedule entire.} The cosmic background does not become more
dangerous when the schedule fails; it becomes the only thing that
happens, which is the sentence the broader Golden program writes
about a biosphere whose muffling agent has been removed.

That is not a coincidence of phrasing. \emph{An incomplete schedule and
the extinctionist's withdrawn world are the same configuration.} Both
are $x$-cannot-be-held; both sit at bare $B_t$; both are
the broader Golden program's case; and both lie in
$\mu G\setminus\mathrm{Base}$ by the downstream Route World hypothesis --- departed
by configuration rather than by event, with no hazard yet fired. The
channel therefore does not merely add a term. It exhibits the model's
apparently smallest channel as conditional on its largest improbability.
\end{remark}

\subsection{Failure incurs debt, not weight}\label{sch:subsub:debt}

What a wrong gate costs can now be said, and the natural first answer is
the wrong one. Treating a wrong gate as multiplying an onset weight
makes that weight a primitive, a number the model is handed rather
than one it produces, and leaves the cascade of
the downstream Route World hypothesis inexpressible, since a single multiplicative
factor has nowhere to put the recursion by which an Imperfection
Magnitude Drop's residual carries continuations of its own.

\begin{definition}[Debt; \textbf{[S]}]\label{sch:def:debt}
A wrong gate does not suppress a weight. It \emph{incurs debt}: $r$
further gate-successes are owed, from the state now occupied, before
completion is reached. Write $Z_t$ for the outstanding debt, with
$Z_0=n$. Each firing consumes one gate and, with probability $1-\pi$,
adds $r$:
\begin{equation}\label{sch:eq:debt}
Z\;\longmapsto\;
\begin{cases}
Z-1 & \text{with probability }\pi,\\
Z-1+r & \text{with probability }1-\pi .
\end{cases}
\end{equation}
\emph{Completion} is the event $Z=0$.
\end{definition}

The natural-historical reading is the one that motivates it. Human
extinction does not leave the biosphere holding a fractional prospect;
it leaves it owing a further sequence, some other lineage must run
its own improbable gymnastics from wherever the world now stands, and
the onset state recedes by exactly that sequence's length. Debt
postpones a deadline. It does not discount a prize.

$r$ need not be a primitive, and the model is better if it is not. Read
the schedule as a tree whose paths each carry a definite number of
events, differing between paths: a wrong turn then places the biosphere
on a longer path, and \emph{the excess length is the debt}. On that
reading \eqref{sch:eq:debt} is the induced dynamics on \emph{events
remaining} rather than a further posit, $r$ is a structural property of
the tree rather than a parameter, and whether it correlates with the
rates, assumed away above, becomes a question about how the tree
is built. Nothing downstream changes, since the induced dynamics is
\eqref{sch:eq:debt}.

Two things follow immediately and neither was available before. \emph{The win condition is binary.} A branch either clears its debt or
does not, and any branch that clears it completes --- however deep the
debt ran, and with no residual weight to carry. The grading that a
weight formulation would posit emerges instead from the
\emph{distribution} over branches. \emph{And the recursion is structural.} Debt spawns debt: the extra
gates are gates, they can go wrong, and going wrong incurs more. That is
the cascade's shape, and the object it makes is standard.

\begin{proposition}[The debt process is Galton--Watson; \textbf{[P]}]
\label{sch:prop:gw}
Under Definition~\ref{sch:def:debt}, $Z$ is the population of a
Galton--Watson process with offspring probability generating function
\begin{equation}\label{sch:eq:pgf}
f(u)\;=\;\pi+(1-\pi)\,u^{\,r},
\qquad\text{mean offspring}\quad \varkappa\;=\;r(1-\pi).
\end{equation}
Completion is extinction of $Z$. Hence completion occurs almost surely
if and only if $\varkappa\le1$, and for $\varkappa>1$ the probability of never
completing is $1-q^{\,n}$, where $q$ is the least root of $q=f(q)$ in
$[0,1]$.
\end{proposition}

\begin{remark}[A phase transition in the schedule's own structure]
\label{sch:rem:criticality}
$\varkappa=r(1-\pi)=1$ is a sharp threshold, and it is not a restatement of
anything already in the paper. Below it, debt clears with probability one: however unlucky a biosphere
is, it owes a finite sequence and eventually pays it. Above it, a
positive fraction of branches accumulate debt without bound, cycling
into ever-deeper obligation, each new gate as fallible as the last, and those branches never reach the state at which $x$ can be held. The
transition sits at $r=1/(1-\pi)$: a biosphere may tolerate large debt
per failure if failures are rare, or frequent failures if each is cheap,
and the product is what decides.

The threshold is structurally the anatomy threshold's sibling. $\rho=0$
divides configurations by whether acquired capacity protects; $\varkappa=1$
divides them by whether incurred debt is payable. Both are sign tests on
a structural parameter, both are sharp, and neither is visible in any
magnitude the system displays.
\end{remark}

\begin{remark}[Under log-spread rates the criterion is the geometric
mean; \textbf{[P]}]\label{sch:rem:bpre}
Definition~\ref{sch:def:rates} redraws the rates at every gate, so
$\pi_k$, and with it $\varkappa_k=r(1-\pi_k)$, is not a deterministic
sequence but an i.i.d.\ random one. The debt process is then a branching
process in a \emph{random} environment rather than a varying one, and the
criterion changes accordingly. By the Smith--Wilkinson criterion, under the standard nondegeneracy
and integrability conditions on the environment, completion is almost
sure when
\begin{equation}\label{sch:eq:bpre}
\E\bigl[\log \varkappa\bigr]\;\le\;0,
\end{equation}
the critical case included, and $\E[\log\varkappa]>0$ permits a positive
probability of never completing: the criterion is the \emph{geometric}
mean of the offspring means, and not $\E[\varkappa]<1$.

The two disagree, and the direction of the disagreement matters. By
Jensen $\E[\log \varkappa]\le\log\E[\varkappa]$, so a schedule may complete with
probability one while its \emph{arithmetic} mean offspring exceeds
unity. Using $\E[\varkappa]$ is therefore not a conservative simplification but
a wrong answer in the pessimistic direction: it condemns schedules that
in fact clear.\footnote{Simulation, Poisson offspring, $4000$ trials
each. With $\varkappa$ equal to $0.1$ or $5$ with equal probability,
$\E[\varkappa]=2.55$ and $\E[\log \varkappa]=-0.348$: the arithmetic criterion calls it
supercritical, the geometric criterion subcritical, and completion is
observed with frequency $1.000$. Where the two criteria agree, $\varkappa\in\{0.5,3\}$ and $\varkappa\in\{0.9,1.6\}$, both geometrically supercritical, completion frequencies are $0.738$ and $0.690$.}

One caution about what the criterion certifies. \eqref{sch:eq:bpre}
settles whether the debt clears \emph{eventually}; it does not say when.
Under a random environment the probability of not having completed by
stage $n$ decays at the tilted rate
$\gamma^\ast=-\min_{0\le s\le1}\log\E[\varkappa^{\,s}]$, which is bounded
above by $|\E\log\varkappa|$ and can fall far below it: for the
$\varkappa\in\{0.1,5\}$ environment above, $\gamma^\ast=0.0158$ at
$s^\ast=0.092$ against $|\E\log\varkappa|=0.347$, a factor of
twenty-two.\footnote{The rate identity is the standard large-deviation
form for subcritical branching in a random environment; the polynomial
correction depends on the subcriticality class.} Certain completion and
timely completion are therefore different questions, and only the first
is settled here \textbf{[O]}.

The reading is the one this work keeps arriving at. A geometric mean is
what a typical realization experiences; an arithmetic mean is what an
ensemble average reports. The quenched quantity governs, exactly as
Remark~\ref{sch:rem:gap2} says it does for the completion weight, and
here it governs the dichotomy itself rather than only the rate.
\end{remark}

\paragraph{Unrecoverable failures, and a doom weight from the
schedule alone.}\label{sch:rem:defective}
Definition~\ref{sch:def:debt} makes every failure recoverable at a
price. Some plausibly are not, an event arriving before its
prerequisites may foreclose rather than cost, and the natural
encoding is $r=\infty$, which is to say a \emph{defective} offspring
law. Let a gate resolve correctly with probability $\pi$, recoverably
with probability $q_1$ at debt $r$, and unrecoverably with probability
$q_2=1-\pi-q_1$. Then
\begin{equation}\label{sch:eq:defective}
f(u)=\pi+q_1u^{\,r},\qquad f(1)=\pi+q_1<1,
\end{equation}
so $u=1$ is not a fixed point and the least root of $u=f(u)$ lies
strictly below one \emph{whatever} the recoverable branch does.
Completion probability is bounded away from one by $q_2$ alone, and
subcriticality of the recoverable branch does not rescue it.

Two consequences. The schedule now produces a positive doom weight
unaided, a nonzero probability of never reaching the onset state,
computed from $(\pi,q_1,q_2,r)$ and nothing else, which is a second
concrete instance of the graded object
Remark~\ref{sch:rem:mvalued} discusses. And the question whether
foreclosing options dominate acquires two answers rather than one: at
fixed mean offspring the second factorial moment enters the survival
bounds, so they may dominate the \emph{probability} of non-completion
while leaving the mean untouched.\footnote{Checked
against direct simulation of \eqref{sch:eq:debt} with a killing branch,
$4\times10^{4}$ trials each: at $(\pi,q_1,r)=(0.85,0.10,3)$ the least
root is $0.9306$ against $0.9310$ simulated; at $(0.70,0.25,2)$,
$0.9046$ against $0.9040$; at $(0.60,0.35,2)$, $0.8571$ against
$0.8557$.}

\subsection{The offset weighting, emergent}\label{sch:subsub:emergent}

The constraint on this formulation was that it reproduce the companion's
offset weighting rather than replace it. It does, and the derivation
also says what the offset is \emph{made of}. By Proposition~\ref{sch:prop:gating} a biosphere pays $B+\Gamma$ while
incomplete --- the background unshielded, and the internal channel at its
bare coefficient. Debt is therefore paid in hazard: clearing $Z$ gates at
rate $\lambda$ costs $D/\lambda$ in expected exposure, and since each
gate's waiting time is exponential the survival weight it contributes
is the Laplace transform $\E[e^{-(B+\Gamma)T}]=\lambda/(\lambda+B+\Gamma)$,
not $e^{-(B+\Gamma)/\lambda}$; the two agree to first order in
$(B+\Gamma)/\lambda$ and differ beyond it.

\begin{proposition}[Completion weight as a least fixed point;
\textbf{[P]}]\label{sch:prop:fixedpoint}
Let $s:=\lambda/(\lambda+B+\Gamma)$ be the survival factor per gate, the
Laplace transform of the gate's exponential waiting time at the
incomplete state's hazard. The expected completion weight from a single
owed gate is the least root $u\in[0,1]$ of
\begin{equation}\label{sch:eq:fixedpoint}
u\;=\;s\,f(u)\;=\;\frac{\lambda}{\lambda+B+\Gamma}\bigl(\pi+(1-\pi)u^{\,r}\bigr),
\end{equation}
and from $n$ initial gates it is $u^{\,n}$. Where the rates are drawn
per gate the fixed point is taken in expectation over the draw,
\begin{equation}\label{sch:eq:fixedpointrandom}
u\;=\;\E\!\left[\frac{\Lambda}{\Lambda+B+\Gamma}\bigl(\pi+(1-\pi)u^{\,r}\bigr)\right],
\qquad\Lambda=\sum_j\lambda_j,\quad\pi=\frac{\lambda_1}{\Lambda},
\end{equation}
so that no single-rate surrogate is needed; to first order in
$(B+\Gamma)/\lambda$ the harmonic mean $(\E[\lambda^{-1}])^{-1}$ is the
surrogate that reproduces the mean log-weight, smaller than the
geometric mean for a lognormal law by the factor $e^{-s_\lambda^2/2}$. For small $B/\lambda$,
\begin{equation}\label{sch:eq:expansion}
-\log u\;=\;\frac{B+\Gamma}{\lambda}\cdot\frac{1}{1-\varkappa}
+O\!\left(\left(\tfrac{B}{\lambda}\right)^{2}\right),
\qquad \varkappa=r(1-\pi)<1 .
\end{equation}
\end{proposition}

\begin{proof}
\eqref{sch:eq:fixedpoint} is the total-progeny functional equation for a
Galton--Watson process weighted by $s$ per individual. For
\eqref{sch:eq:expansion}, write $u=1-\varepsilon$ and expand:
$1-\varepsilon=(1-(B+\Gamma)/\lambda)(1-\varkappa\varepsilon)+O(\varepsilon^{2})$
gives $\varepsilon(1-\varkappa)=(B+\Gamma)/\lambda$. The factor $1/(1-\varkappa)$ is the expected total
progeny per individual, as it must be.
\end{proof}

\paragraph{What the offset is made of, and where the recursion
went.}\label{sch:rem:cefff}
Two effects separate cleanly, and the separation is this subsection's
return.

\emph{Contrast per unit of depth.} A single wrong gate adds $r$ gates,
each exponential at rate $\lambda$, so the added waiting is
$\mathrm{Gamma}(r,\lambda)$ and the weight multiplies by
$\E[e^{-hT}]=(1+h/\lambda)^{-r}$ with $h:=B+\Gamma$. Matching to the
companion's $10^{-c_{\mathrm{eff}}}$,
\begin{equation}\label{sch:eq:cefff}
c_{\mathrm{eff}}\;=\;r\,\log_{10}\!\Bigl(1+\frac{B+\Gamma}{\lambda}\Bigr),
\end{equation}
exactly and with no approximation. The deterministic-exposure reading, $r/\lambda$ of waiting at hazard $B+\Gamma$, giving
$(B+\Gamma)r/\lambda\ln10$, is the $\lambda\gg B+\Gamma$ limit of it
and \emph{overstates} the contrast elsewhere, since $\log(1+x)<x$. The
direction of that error is worth naming: it makes the companion's
collapse look easier than it is, so \eqref{sch:eq:cefff} is what should
be used. The High-Contrast Earth
Hypothesis's ``extremely many orders'' is thereby not a posit but a
reading: one gets many orders when the background is high, the debt
large, or the tempo slow.

\emph{Depth accumulated.} The recursion does not change
\eqref{sch:eq:cefff}; the branching factor $1/(1-\varkappa)$ cancels out of the
per-depth contrast exactly. What it changes is how much depth there is:
expected total gates fired per owed gate is $1/(1-\varkappa)$, so a biosphere
near criticality accumulates depth without the contrast per unit of it
worsening at all. The recursion inflates the exposure, not the exchange
rate --- and inflates it without bound as $\varkappa\to1$.

The separation identifies two independent levers. A biosphere may be
badly placed because each failure is expensive ($c_{\mathrm{eff}}$
large) or because failures compound ($\varkappa$ near one), and the two are not
the same defect.

\begin{remark}[Contrast wants slow gates; survival wants fast ones]
\label{sch:rem:tension}
\eqref{sch:eq:cefff} has two limits and they pull opposite ways, which
is a structural fact about the channel rather than a defect of it. Where gates are fast against the standing hazard, $\lambda\gg B+\Gamma$,
the contrast is $c_{\mathrm{eff}}\approx r(B+\Gamma)/\lambda\ln10$ ---
small. Wrong turns cost little, passage is cheap, and alternatives are
barely suppressed: there is nothing to collapse a sum over.

Where gates are slow, $\lambda\ll B+\Gamma$, the contrast is
$c_{\mathrm{eff}}\approx r\log_{10}\bigl((B+\Gamma)/\lambda\bigr)$ --- a
\emph{log-rate gap}, hence many orders of magnitude whenever the gap is
large. But the same slowness prices survival at
$(\lambda/(B+\Gamma))^{\,r}$ per wrong turn, which is astronomically
small. So the regime supplying the companion's contrast is the regime in which
the realized history had to be enormously lucky, and the two are one
fact seen from either end. That is not a tension to be resolved but the
Miracle phenomenology in rate coordinates: high contrast between natural
histories and low probability of the realized one are not independent
observations.
\end{remark}

\begin{remark}[The collapse is a claim about shares, not about
likelihood]\label{sch:rem:absolute}
A companion point, and the one most easily misread. The sum-collapse
establishes that the realized ordering carries a share of order unity
\emph{of the total}. It says nothing whatever about the size of that
total, and under a log-scale randomization the total is minute. Two facts about extremes make this precise. The best of $m$ draws from a
spread log-rate law is not far above the field: the maximum of $m$
normals sits $\sqrt{2\ln m}$ deviations above the mean asymptotically,
and only $1.54$ deviations at $m=10$ --- so the optimum is not
qualitatively better than a typical option, it is the same law read at
its edge. And the optimum still faces the standing hazard: with
$\lambda\ll B+\Gamma$ the weight of even the best option at a gate is
$\approx\lambda/(B+\Gamma)$, orders below one, and $n$ gates multiply
those orders.

A worked instance makes the two magnitudes visible together. Take
$\log_{10}\lambda$ normal with mean $15$ decades below $B+\Gamma$ and
deviation $3$, with $m=10$, $n=20$, $r=3$. The best option at a gate
sits $4.6$ decades above typical; the weight of the all-best path is
about $10^{-208}$; a single wrong turn costs $c_{\mathrm{eff}}\approx31$
decades; and the collapse condition $c_{\mathrm{eff}}>\log_{10}(n(m-1))
=2.26$ is met roughly fourteen times over.\footnote{Direct simulation,
$2\times10^{5}$ draws; the extreme-value theory agrees to two decimals,
$\mu+1.5388s=-10.38$ against $-10.39$ observed.}

So relative dominance is overwhelming and absolute weight is
$10^{-208}$, and the two are not in tension --- both follow from the
same spread, which is what makes one alternative dominate its rivals and
what makes all of them improbable. The reading to avoid is therefore
explicit: \emph{the highest-weight ordering is not the likely ordering}.
Nothing in the companion's apparatus, and nothing here, converts a share
of a small number into a large one. This is the same discipline
the broader Golden program applies to the survivor's own position, and
it applies with more force here because the numbers are worse.
\end{remark}

\begin{remark}[Which contrast the collapse consumes]\label{sch:rem:n2}
A random contrast does not enter the companion's sum-collapse the way a
deterministic one does, and the difference has to be recorded. An alternative at depth $\Delta$ is suppressed by
$10^{-\sum_{i\le\Delta}c_i}$, so with independent draws the per-tier
factor is $(\E[10^{-c}])^{\Delta}$ and the binomial closes at
$\bigl(1+(b-1)\E[10^{-c}]\bigr)^{n}$. The operative quantity is
therefore $-\log_{10}\E[10^{-c}]$, not $\E[c]$, and by Jensen the first
is strictly smaller. The gap is severe rather than technical: for $c$
uniform on $(0,2M)$ the averaged contrast is $\log_{10}(2M\ln10)$,
growing only \emph{logarithmically} in the spread; for $c$ normal with
mean $\mu$ and deviation $s$ it is $\mu-s^{2}\ln10/2$, negative once
$s\gtrsim\sqrt{2\mu/\ln10}$.\footnote{At $\E[c]=10$: uniform on
$(0,20)$ gives $1.66$; normal with $s=2$ gives $5.45$; normal with $s=3$
gives $-0.36$. Direct sampling understates the collapse, four million
draws return $1.25$ for the last, because $\E[10^{-c}]$ is carried by
a tail sampling does not reach, which is itself the point.}

The resolution is the one this work keeps reaching. The averaged
contrast is dominated by alternatives no one encounters; the claim the
collapse is used for concerns the share carried by the \emph{realized}
configuration, which is quenched. One refinement to ``runs on $\E[c]$'':
the quenched tier weight is a sum of order $\bigl(n(b-1)\bigr)^{\Delta}$
lognormal terms and is extreme-value dominated at small $\Delta$, so the
binding constraint sits at the shallowest tier and carries a fluctuation
allowance,
\begin{equation}\label{sch:eq:quenchedcrit}
\bar c\;>\;\log_{10}\bigl(n(b-1)\bigr)
\;+\;\frac{s\sqrt{2\ln\bigl(n(b-1)\bigr)}}{\ln10},
\end{equation}
relaxing toward the $\E[c]$ form as $\Delta$ grows, since the fluctuation
scales as $\sqrt\Delta$ against a mean linear in $\Delta$. At $s=3$,
$n=20$, $b=10$ the threshold moves from $2.26$ to $6.45$, which
$\E[c]=10$ still clears. The companion's
own hedge, that its deterministic use is the union-bounded envelope
of an almost-sure claim, points the same way. So the collapse
survives randomization in the quenched reading and fails in the annealed
one, and a reader wanting the annealed statement needs a contrast far
larger than $\E[c]$ suggests \textbf{[O]}.
\end{remark}

\begin{remark}[This is the graded cascade, in one concrete case]
\label{sch:rem:mvalued}
the downstream Route World hypothesis keeps the cascade's Boolean recursion and the
$\varepsilon$-grading of the downstream Route World hypothesis in separate
layers, and records what would join them: a fixed point of an
$M$-valued doom operator, called open \textbf{[O]} and named there as
the natural next question.

\eqref{sch:eq:fixedpoint} is an instance of exactly that object. The map
$u\mapsto s\,f(u)$ is monotone on $[0,1]$, its least fixed point is the
graded analog of $\mu G$'s least fixed point on the Boolean lattice,
and the grading it carries is a weight rather than a label. The
completion probability is therefore \emph{computed} as a fixed point
rather than assigned, which is what the two-layer remark asks for.

The claim should be taken at its size. This is a single-type process
carrying one number, not the $M$-valued operator over the full lattice
of states the general construction would need; what is exhibited is the
first concrete case and the shape it confirms. What the case does supply is the general construction's
\emph{definition}, and the choice is not the obvious one. The Boolean
cascade characterizes $\mu G$ modally: every course admits a foreclosing
continuation, which is a fact about the state space. A graded version
cannot be built on that, because a modal fact carries no measure. It can
be built on the dynamical reading, \emph{all paths lead to} the
Departure Point, rather than all courses make it exist, since that is
a hitting statement, and the measure of paths hitting a set is
canonically the \emph{least} fixed point of the associated operator.
\eqref{sch:eq:fixedpoint} is exactly such a hitting measure, which is why
it arrived as a least fixed point without anyone imposing it.

Two things follow for the general construction. The operator to build is
the one whose fixed point is the hitting measure of the departure set,
and Kleene iteration from $\bot$ reaches it wherever the weight algebra's
multiplication is Scott-continuous. And the projection back to the
Boolean layer is a \emph{threshold} rather than a linear map, Boolean
$\mu G$ is where the surviving weight is zero, not where the graded value
takes a particular magnitude, so the consistency requirement is
$\mathrm{eval}(u)=\one\{u=0\}$, and it should be imposed at definition
time rather than proved afterwards. Whether the construction so specified
generalizes to the lattice remains open \textbf{[O]}.
\end{remark}

\begin{remark}[Which rate an observer should use]\label{sch:rem:gap2}
The completion weight $u^{\,n}$ is an expectation over branches, and
near criticality it is dominated by the rare branch that incurred little
debt: by \eqref{sch:eq:expansion} the \emph{mean} total progeny is
$1/(1-\varkappa)$, while the median sits far below it. An observer reasoning
about their own history should therefore not use the mean. Expected
completion weight is carried by branches that got lucky early, and a
biosphere is not entitled to assume it is on one --- the same
quenched--annealed discipline the reference sheet records for the
survivorship transform.\footnote{The same gap is recorded for the survivorship transform in the broader Golden program; the two are instances of one phenomenon rather than one result applied twice.}
\end{remark}

\subsection{The composite}\label{sch:subsub:composite}

The schedule does not win Golden-Point probability. It wins the
probability of reaching the one state at which $x$ stops decaying, and
everything the model already proves takes over from there. That gives
the total a two-stage shape, schedule then model, with $\tau$ the
completion time and $X_\tau=(\mu_\tau,D_\tau)$ the \emph{model's} state
at that moment, the Kolmogorov pair of Definition~\ref{imp:model}
and not the schedule's debt $Z$, which is zero at $\tau$ by
construction.\label{sch:eq:composite} The second stage is
Theorem~\ref{imp:twocond} evaluated at the completion state, and the
composite is written as a kernel and not as a product, because the
completion state depends on how the schedule was cleared. Let $K_{\mathrm{sch}}(d\mu,dx,d\rho)$ be the \emph{surviving-completion
subprobability kernel}, the law of the full state the schedule leaves
at $\tau$, driver, strength, and anatomy, weighted by survival to
$\tau$,
$K_{\mathrm{sch}}(A)=\E[e^{-\Lambda_\tau}\one\{\tau<\infty,\,(\mu_\tau,x_\tau,\rho_\tau)\in A\}]$.
Then
\begin{equation}\label{sch:eq:threecond}
\PP\bigl(\text{advance}\bigr)
=\int\one\{\rho>0\}\,V_{\mathrm{drv}}(\mu)\,K_{\mathrm{sch}}(d\mu,dx,d\rho),
\end{equation}
the driver-viability functional taking the driver coordinate $\mu$ as
its argument, the strength coordinate entering through the survival
weight and the prefactors and not through the event-level factor,
with $V_{\mathrm{drv}}$ the driver-viability functional of
Theorem~\ref{imp:twocond}. The three ingredients are then visible without
being multiplied: schedule failure is mass $K_{\mathrm{sch}}$ fails to
deliver, bad anatomy is the indicator, and bad driver dynamics is
$V_{\mathrm{drv}}$; and $u^{\,n}$ of \eqref{sch:eq:fixedpoint} is the
total mass of $K_{\mathrm{sch}}$ in the homogeneous model, a
survival-discounted completion weight, $\E[s^N\one\{\text{completion}\}]$
in the gate-count representation, which reduces to the completion
probability only at $s=1$. The composite becomes a product only under
independence or a deterministic completion state, neither of which the
model supplies.

\paragraph{The three factors are not of one kind.}\label{sch:rem:threefactor}
Anatomy \emph{annihilates}: $\rho\le0$ gives zero, not a small number,
by Theorem~\ref{imp:twocond}. The schedule \emph{attenuates}: a
wrong gate costs debt, and debt cleared is completion, so the factor is
drastically minute and never quite zero --- which is what it is to be an
Imperfection Magnitude Drop in the companion's exact sense \cite{Nexus}.
So the schedule factor is the residual $\varepsilon$ of
the downstream Route World hypothesis computed rather than named, and the
channel's connection to the cascade is an identification rather than an
analogy.

\paragraph{The decomposition is hierarchical, not a product.}\label{sch:rem:hierarchical}
\eqref{sch:eq:threecond} should be read exactly. It is not an
independence factorization, since $X_\tau$ depends on how long the debt
took to clear, but a hierarchical decomposition through the kernel, and
what makes it usable is that the influence runs one way. The schedule's law is autonomous: $\lambda_k$
and $\pi_k$ consult the gate count and nothing else, so the schedule may
be solved first, without reference to the driver. The leak then carries
the schedule's verdict into the driver through $\gamma_{k_t}$, and the
model runs from the state that leaves. Schedule, then driver, then
model, with no arrow back.

That one-directionality is what a state-dependent schedule threatens,
and it threatens it in only one of the two ways such dependence can be
read. Let $\pi$ depend on the schedule's own \emph{configuration}, which gates are closed, in what order, with what debt outstanding, and the schedule remains autonomous, since configuration is generated by
the schedule itself. It becomes a multitype branching process, with
types the configurations and criticality governed by the top Lyapunov
exponent of products of mean matrices rather than by the scalar
criterion of \S\ref{sch:subsub:open}, and the hierarchy survives intact.
Let $\pi$ instead depend on the \emph{driver}, on $x_t$ or $\mu_t$, and the schedule consults what it was supposed to precede. The
hierarchy collapses, \eqref{sch:eq:action} becomes a genuine control
problem, and the optimal history must trade action spent early against
exposure shortened.

The baton-passing picture that motivates state dependence is of the
first kind: what an event does depends on the state its predecessor
left, which is configuration. So the version the natural-historical
reading requires costs a multitype upgrade and not the control problem.
Whether the driver \emph{also} influences gate outcomes is a modeling
question about the world rather than a mathematical one, and the answer
adopted here is that it does not. Gate outcomes are stochastic in the
direct sense: what makes an option the realized one is a fact about the
onward schedule it opens, not about the credit held when it fires. The
hierarchy therefore holds, \eqref{sch:eq:action} does not become a
control problem, and the multitype upgrade is the whole cost of state
dependence.

One refinement that adoption invites is recorded rather than pursued. If
an option is realized in virtue of the ease of the gates it opens, then
correctness is not a label attached to one of the $m$ branches but a
\emph{derived} property of the subtree below it --- which makes the
schedule a tree carrying a value function rather than a sequence with a
marked option, and makes the realized path something survivorship picks
out rather than something given in advance. That is consistent with
Definition~\ref{sch:def:rates}, whose independence claim concerns an
option's \emph{own} rate and not its descendants'. Whether the tree
formulation reproduces the sequence one at the level of
\eqref{sch:eq:fixedpoint} is open \textbf{[O]}.

\begin{remark}[Finite horizons select the schedule; the asymptote
excludes its absence]
\label{sch:rem:asymptote}
The channel prices what climbing costs. Whether anything \emph{selects}
for climbing is a separate question, and the answer separates the two
conditioning events this work uses.

\emph{Finite horizons select but do not exclude.} Condition on
$\{T>t\}$ for any finite $t$; by Proposition~\ref{sch:prop:threshold}
the conditioning favors completion increasingly with $t$, but never
removes the alternative from the support. A biosphere that never completes its schedule sits at bare
background with its internal channel at its bare coefficient, by
Proposition~\ref{sch:prop:gating}, and survives to $t$ with probability
$e^{-(B+\Gamma)t}>0$. The un-climbed branches therefore carry
positive weight under every finite-horizon conditioning, and the lucky
route, survive by not being struck, and never pay the schedule at all, remains in the support throughout.

\emph{The asymptote excludes.} Condition instead on the viability event
$\Imm=\{\Lambda_\infty<\infty\}$ of Definition~\ref{imp:model}. By
Proposition~\ref{sch:prop:gating},
$\PP\bigl(\Imm\mid\text{schedule incomplete}\bigr)=0$ \emph{exactly} ---
not small. The exclusion is at the level of support rather than of
weight, and it is the only conditioning in this work that achieves it. So the schedule is selected for by the asymptotic condition and by
nothing weaker. Between the finite horizon and the asymptote there is a
threshold, and it is the mechanism by which conditioning on survival
pushes a lineage up the schedule at all.

\begin{proposition}[The survival-time threshold for pushing the
schedule; \textbf{[P]} given the model]\label{sch:prop:threshold}
Consider a biosphere whose histories fall into two classes: the
\emph{unpushed}, which never complete the schedule and survive to $t$
with probability $e^{-(B+\Gamma)t}$, and the \emph{pushed}, which
complete it, an event of unconditional probability $p$, and thereafter
survive with probability $\asymp Ct^{-1/4}$ in the critical regime
(the persistence tail of \cite{SurvA}) or with positive probability in the viable
one. Then:
\begin{enumerate}[label=\emph{(\roman*)},itemsep=2pt]
\item the posterior weight of the pushed class given survival to $t$ is
\begin{equation}\label{sch:eq:pushposterior}
\PP(\text{pushed}\mid T>t)=\frac{p\,Ct^{-1/4}}{p\,Ct^{-1/4}+(1-p)e^{-(B+\Gamma)t}}
\;\longrightarrow\;1\qquad(t\to\infty);
\end{equation}
\item there is a threshold $T_\ast$, the solution of
$p\,CT_\ast^{-1/4}=(1-p)e^{-(B+\Gamma)T_\ast}$, so that
\begin{equation}\label{sch:eq:threshold}
T_\ast=\frac{1}{B+\Gamma}\Bigl[\log\frac{1-p}{pC}+\tfrac14\log T_\ast\Bigr]
\;=\;\frac{\log(1/p)}{B+\Gamma}\,\bigl(1+o(1)\bigr)\qquad(p\to0),
\end{equation}
below which the conditioned law favors the unpushed class and above
which it favors the pushed class, the switch being sharp because one
weight is exponential in $t$ and the other polynomial: rarity of
pushing, however extreme, enters the threshold only through
$\log(1/p)$, so a one-in-a-trillion push moves $T_\ast$ by about
$28/(B+\Gamma)$ and no more;
\item within the pushed class, conditioning on survival to $t$ weights
each route $r$ through the schedule by its survival weight, and as
$t\to\infty$ the ratio of two routes' weights converges to the ratio of
their persistence-harmonic prefactors at completion,
$\hm(X^{(r)}_\tau)/\hm(X^{(r')}_{\tau'})$ (the conditioned-law theorem of \cite{SurvA}), so the
conditioned law \emph{tilts} the completed-route distribution by the
persistence-harmonic prefactor at completion, toward better completion
states; two routes with prefactors in the ratio $2{:}1$ keep limiting
weights in that ratio, and concentration on the route maximizing
\eqref{sch:eq:action}, the optimal route, follows only in a further
limit in which the prefactor contrasts diverge, entered as
\textbf{[S]};
\item conditioning on asymptotic survival, $t\to\infty$, lies past
$T_\ast$ for every parameter value, so the asymptotically conditioned
law prefers the pushed class outright, tilted toward better completion
states, with the unpushed class excluded at the level of support by
Proposition~\ref{sch:prop:gating}.
\end{enumerate}
\end{proposition}

\begin{proof}
(i) is Bayes on the two-class mixture with the stated survival
probabilities; the limit follows because $t^{-1/4}$ dominates
$e^{-(B+\Gamma)t}$. (ii) is the balance equation of (i) solved for
$T_\ast$, with the asymptotic form from neglecting the logarithmic
term. (iii): conditioning on $\{T>t\}$ weights a route by
$\PP(T>t\mid\text{route})$, which by the conditioned-law theorem of \cite{SurvA} is
$\asymp C\,\hm(X_\tau)\,t^{-1/4}$ with the prefactor set by the completion
state, so ratios of route weights converge to ratios of $\hm$ at
completion, which is the tilt; concentration would require the ratio of
the maximal prefactor to every other to diverge, which no hypothesis
here supplies. (iv): $T_\ast<\infty$ for every $p>0$, $C>0$,
$B+\Gamma>0$, and the support exclusion is
Proposition~\ref{sch:prop:gating}.
\end{proof}

The proposition is the mechanism behind the ladder. A lineage
conditioned on asymptotic survival is one whose schedule was pushed,
since every unpushed history is exponentially suppressed past $T_\ast$
and excluded at the asymptote; and among pushed histories the
conditioning tilts toward the better-completing routes, since the
survival weight of a route is its harmonic prefactor at completion,
and reaches the optimal route only in the limit of diverging
contrasts. Below the
threshold, the same conditioning says the opposite, that not pushing is
the likelier history, because pushing is rare and the exponential
penalty for not having pushed has not yet accumulated. This is why an
amoeba lineage, read under asymptotic conditioning, is read as having
been pushed up the ladder, and why the situational contingencies the
Optimal-Earth companion catalogues fall out of the conditioning rather
than being posited beside it: reaching the threshold at all is
exponentially less likely on every unpushed route. The crossover
recedes as $C$ falls, and $C$ is set by the state completion leaves
(Remark~\ref{sch:rem:initialdata}), so a biosphere completing into a
poor state finds the schedule worth having climbed only at horizons
that run toward the asymptote; the threshold's dependence on $C$ is
again logarithmic.

The reading is the one that matters for what this channel is for.
\emph{Nothing about surviving a finite stretch gives a reason to climb.}
Only conditioning on there being no last moment makes the climb pay ---
which is why the schedule appears in the record of a biosphere whose
history was conditioned, and would not appear in the record of one that
had merely lasted.
\end{remark}

\begin{remark}[The schedule supplies the initial data the model assumes]
\label{sch:rem:initialdata}
\eqref{sch:eq:composite} makes an identification that should be stated plainly:
\emph{the completion state is the initial datum} that
the survivorship companion \cite{SurvA} takes as given. The model begins where the
schedule ends, and the schedule delivers not a point but a distribution
over points --- shaped by how long the debt took to clear and how far
$x$ decayed while it did.

By the survivorship companion \cite{SurvA} the initial data are gauge: they act
through prefactors and timescales, never through exponents. So the
schedule cannot move the persistence exponent $\tfrac14$, cannot move
the trichotomy, and cannot move the two-condition theorem. What it sets
is the prefactor --- and by the survivorship companion \cite{SurvA} the prefactor
is proportional to the persistence-harmonic function $\hm$ of the
starting state. That is the same $\hm$ the Golden Score computes
(the broader Golden program), and the broader Golden program records as
its principal cost that no map exists from a decision to a point in
$(\mu,D)$. The schedule produces exactly such points. Whether decisions
bearing on a biosphere's schedule, on $\lambda$, on $\pi$, on debt
incurred, thereby acquire computable harmonic values is the first
place a candidate identification map might be found. We flag it as a
lead and claim nothing \textbf{[O]}.
\end{remark}

\begin{remark}[The variational form, stated and not solved]
\label{sch:rem:variational}
The second factor of \eqref{sch:eq:composite} is a conditioned path
problem, and the survivorship companion \cite{SurvA} already names the tool: the
adverse-drift rate and its front-load-then-coast profile come from the
Cameron--Martin action by Cauchy--Schwarz, the model's analog of an
Euler--Lagrange computation and the natural processing step wherever a
conditioned path law is wanted at large-deviation scale.

The composite asks for that computation across both stages. One seeks
the path minimizing
\begin{equation}\label{sch:eq:action}
S[\mu]\;=\;\int_0^{\infty}\left[\frac{\dot\mu_t^{2}}{2\sigma^{2}}
\;+\;h\bigl(x_t,k_t\bigr)\right]dt
\end{equation}
subject to \eqref{sch:eq:leak}, to the debt dynamics
\eqref{sch:eq:debt}, and to $\Lambda_\infty<\infty$: a least-action path
through a landscape whose running cost is the hazard and whose
constraint is the schedule. The pathfinding reading is exact rather than
figurative --- the action is a length and the optimal history a geodesic
in it.

Two features of the answer are predictable without solving it, and both
are checks on a solution's behavior. The running cost during
incompleteness is $B$ and is \emph{independent of the driver}, since $x$
is not held and $\varsigma_b(0)=1$; so the schedule stage contributes a
term the driver cannot optimize against, and the variational freedom
lives entirely after $\tau$. And because \eqref{sch:eq:composite}
factorizes, the optimal post-completion profile is the existing one,
started from the distribution Remark~\ref{sch:rem:initialdata}
describes rather than from a chosen point.

Whether the stages remain separable once $\lambda$ or $\pi$ depends on
the state, which the baton-passing picture requires, is the
substantive question and is open \textbf{[O]}. If they do not,
\eqref{sch:eq:action} becomes a genuine control problem rather than a
factorized one, and the optimal history would trade early action against
schedule exposure in a way the present separation forbids.
\end{remark}

\subsection{The internal channel is probably a schedule too}
\label{sch:subsub:internalschedule}

One structural claim precedes the open list, because it is
the channel's clearest suggestion about the rest of the model and because
it is an inference the Continuity Principle licenses rather than a
speculation laid beside it. the downstream Route World hypothesis holds that hazard-lowering across
mediations is one activity and not a family of analogs. The schedule
represents pre-onset hazard as gates. If continuity is taken at the width
the downstream Route World hypothesis states it, then post-onset internal
hazard should be representable the same way, and the trough's smooth
exponential should be a coarse-graining of gates not yet identified.

\begin{remark}[The trough's exponential is geometric compounding;
\textbf{[S]}]\label{sch:rem:internalgates}
Take $h^{\mathrm{int}}$ proportional to the count of \emph{unpatched
failure modes} carried at capacity $x$. Suppose each increment of
capacity both creates modes and passes gates that patch them, and write
$q$ for the fraction left unpatched per unit of capacity. Then
$\text{unpatched}(x)=\text{unpatched}(0)\,q^{\,x}$, so
\begin{equation}\label{sch:eq:internalgates}
h^{\mathrm{int}}(x)\;=\;\Gamma q^{\,x}\;=\;\Gamma e^{-\rho x},
\qquad \rho=\ln(1/q).
\end{equation}
So geometric compounding of per-gate patch rates \emph{would} produce
the exponential of Definition~\ref{imp:model}, with $\Gamma$ the count at
$x=0$. The construction is offered at that strength and no more: it
exhibits a mechanism yielding the observed form, not evidence that this
is the mechanism. Definition~\ref{imp:model} remains the paper's statement of
the internal channel, nothing downstream is rewritten in gate
coordinates, and the construction's use is heuristic --- it says what one
would look for, and becomes a derivation only if the gates are named.

Three things follow, and the second is the reason to record this. \emph{The anatomy trichotomy becomes a compounding test.}
the survivorship companion \cite{SurvA} divides on the sign of $\rho$;
under \eqref{sch:eq:internalgates} that is $q$ against $1$ --- patching
outpacing creation, matching it, or falling behind. \emph{And it is then the same test the schedule already runs.}
Remark~\ref{sch:rem:criticality} calls $\varkappa=1$ the anatomy threshold's
sibling. Under this construction they are not siblings but one test
applied twice: both ask whether a compounding factor exceeds unity, and
they differ only in what indexes the compounding --- gate count for the
schedule's debt, capacity for the trough's unpatched modes. That is what
continuity predicts, and finding it is a check on the principle rather
than a use of it.

\emph{The construction agrees with the one the paper already has.}
the broader Golden program derives $\rho=-(\alpha-\beta)$ from boundary
scaling --- an extractive domain growing as $e^{\alpha x}$ against
enforcement growing as $e^{\beta x}$. Under \eqref{sch:eq:internalgates}
that is $q=e^{\alpha-\beta}$, so $q>1$ exactly when $\alpha>\beta$, and
the two microfoundations return the same condition by different routes.
Neither is derived from the other.
\end{remark}

What is missing is the thing the reformulation is for: \emph{which
events}. The schedule channel names its gates as natural-historical
transitions, and the internal channel's gates would be whatever a
biosphere must get right for acquired capacity not to be turned inward
--- institutional, presumably, and the broader Golden program is where a
candidate list would come from. Until they are named, the smooth rate is the correct representation, and the reason needs precise statement: a
smooth hazard is the \emph{annealed} description of an unidentified jump
process, so by the discipline of Remark~\ref{sch:rem:gap2} it flatters
the typical history. An observer using Definition~\ref{imp:model}'s trough to
reason about their own biosphere is using a mean carried by branches that
patched luckily \textbf{[O]}.

\subsection{What is open}\label{sch:subsub:open}

\emph{Depth-dependent debt.} $r$ and $\pi$ are constant, so all failures
cost alike. Two quantities must be kept apart before that is repaired,
since they are easily run together and only one of them is ordered by
position. \emph{Stakes}, what failing forfeits, rise monotonically
along the chain, by the broader Golden program: a late step carries
the improbability of every step behind it. \emph{Debt}, what
recovering costs, does not. An early gate may carry enormous $r$: a
biosphere without its moon, or with one a third the size, owes a
compensating sequence for axial stabilization and tides at the
\emph{first} gate rather than the last, and it is not obvious that any
finite sequence discharges it. Severity is not ordered by position, and
a model that assumed it was would misplace the worst gates.

Repairing this makes $Z$ a branching process in a varying environment,
$\varkappa_k=r_k(1-\pi_k)$ a sequence rather than a number, and the criterion a
condition on \emph{products} rather than on any single term: writing
$S_k=\prod_{j<k}\varkappa_j$, completion is almost sure when
$\sum_k S_k^{-1}=\infty$ and has positive failure probability when the
sum converges, under a uniform second-moment condition on the offspring
laws. For constant $\varkappa$ this recovers $\varkappa\le1$. The reading is worth
having: one catastrophic gate does not doom a schedule whose remaining
gates are benign enough, since it is the geometric mean of the $\varkappa_k$
that decides, a single spike is survivable in a way a raised floor is
not, with the caveat that under Definition~\ref{sch:def:rates} the
environment is random rather than merely varying, so the operative
criterion is \eqref{sch:eq:bpre} and the sum below describes the
deterministic case. The question of whether a deeper drop costs more
\emph{per unit of debt} is answered there rather than here, and answered
negatively: the rates are redrawn at each gate and correctness is
independent of them, so severity is neither constant nor monotone in
depth but resampled, and it is the distribution rather than any trend
that carries the structure. The same harmonic sum gives the rate and not
merely the dichotomy:
two-sided bounds of Agresti type put the probability of not having
completed by stage $n$ comparable to $\bigl(\sum_{k\le n}S_k^{-1}
\bigr)^{-1}$, up to the second-moment factor, so the quantitative form
comes free with the criterion and is what a bridge to
Proposition~\ref{sch:prop:exponent}'s timing statements would consume
\textbf{[O]}.

\emph{State-dependent gates.} $\lambda(\cdot)$ and $\pi(\cdot)$ are what
the baton-passing picture requires and what
Remark~\ref{sch:rem:variational} identifies as deciding whether the
composite factorizes. \emph{The leak is assigned rather than emergent}, and the incomplete
regime's asymptotics remain unknown (Remark~\ref{sch:rem:limit}). \emph{And the parameters are unmeasured.} $B$, $\lambda$, $\pi$, $r$ and
$n$ are what the companion's serial-attribute program
(the downstream Route World hypothesis) would supply --- and \eqref{sch:eq:cefff}
makes that program's task easier rather than harder, since
$c_{\mathrm{eff}}$ need no longer be estimated directly but follows from
three quantities a historian might plausibly bound.
the survivorship companion \cite{SurvA}'s reservation applies with full force
throughout: a scored schedule index would be a measurement apparatus
optimized against, which is the failure the channel exists to price.


\subsection{The modern-era term, formalized}\label{sch:subsub:modernterm}

The three passages that follow this subsubsection carry the
schedule's modern-era argument in prose --- the live-term
hypothesis, the identification of accordance-legible golden
causality with the term's operation, and the assignment of the
term's last leg to the technoanatomy --- and a reader is owed the
formal content underneath them before the prose is read. The
argument of the subsection as a whole then runs in one line: the
schedule is a sequence of gated clocks (Definition~\ref{sch:def:schedule})
with log-spread rates (Definition~\ref{sch:def:rates}) whose
incompleteness leaks the credit (Definition~\ref{sch:def:leak}) and
gates the external channel (Proposition~\ref{sch:prop:gating}); the
term's contribution to the hazard of Departure is the derailment
hazard defined next; the hypothesis that the schedule is not yet
exhausted is entered on evidence rather than principle; under it the
remaining gates act on the anatomy rather than the credit, with a
settlement law the conditioning front-loads; and none of this
disturbs the ceiling, whose residual the last passage classifies.

\begin{definition}[Derailment hazard, and anatomy-directed gates;
\textbf{[S]}]\label{sch:def:derailment}
With the schedule in the state $k_t$ of Definition~\ref{sch:def:schedule},
the \emph{derailment hazard} is the instantaneous rate at which any
wrong option at the open gate fires first,
\begin{equation}\label{sch:eq:derailment}
\sigma_t\;:=\;\lambda_{k_t}\bigl(1-\pi_{k_t}\bigr),
\end{equation}
and it is the schedule's contribution to the hazard \emph{of the
Departure Point} in the implication sense of
Definition~\ref{imp:goldenpoint}: a derailment is a Departure-sufficient
occurrence carrying the Aethic weight of the derailed branch, which
under the extra-successes reading adds gates to the schedule rather
than closing it outright. The term is high exactly when a gate with a
large total rate and a small favorable share stands open, and that is
the formal content of the phrase ``high stakes toward landing the next
Miracle.'' Gates $k\le n$ are \emph{credit-directed}: their closure
lowers the leak of \eqref{sch:eq:geomleak}. Gates $k>n$, if any, are
\emph{anatomy-directed}: their resolution does not act on the credit
but on the anatomy, the option $j$ that fires at gate $k$ contributing
an increment $\Delta_{k}^{(j)}$ with $\Delta_{k}^{(1)}>0$ for the
realized-favorable option and the wrong options adverse, so that
\begin{equation}\label{sch:eq:rhoprocess}
\rho_t\;=\;\rho_0+\sum_{n<k\le k_t}\Delta_{k}^{(j_k)},
\qquad
\rho_\infty\;=\;\lim_{t\to\infty}\rho_t,
\end{equation}
$j_k$ being the option that fired at gate $k$. The two-condition
theorem reads on $\rho_\infty$.
\end{definition}

\begin{proposition}[Front-loaded settlement under monotone coupling;
\textbf{[C]} given the stated conditions]\label{sch:prop:settlement}
Under Definition~\ref{sch:def:derailment} and the live-term
hypothesis, suppose that the anatomy increments at the anatomy-directed
gates are totally ordered, that the law of each subsequent gate is
conditionally independent of the past given the current anatomy, and
that the cumulative hazard is monotone decreasing in every favorable
increment and in every earlier favorable resolution time. Then
conditioning on survival tilts the law of the final anatomy upward and
the law of the crossing time to $\rho>0$ earlier in the sense of
first-order stochastic dominance. If in addition the survival weight
is log-supermodular, totally positive of order two, in the increments
and the resolution times jointly, the tilt holds in the
likelihood-ratio order. Coordinatewise monotonicity with conditional
independence gives the dominance and not the likelihood-ratio order;
the stronger conclusion needs the stronger hypothesis. Without the
conditions the one-gate comparison below still holds, but its
iteration over gates is not a theorem.
\end{proposition}

\begin{proof}[Sketch]
Two paths identical except at one anatomy-directed gate differ in
accumulated hazard by the integral of
$\Gamma(e^{-\rho^{\mathrm{adv}}x_t}-e^{-\rho^{\mathrm{fav}}x_t})$ over the
time the two anatomies differ, which is positive where $x_t>0$ and
negative where $x_t<0$; the proposition's hypothesis that the
cumulative hazard is monotone decreasing in every favorable increment
is exactly the requirement that the integral be positive, which holds
when $x_t\ge0$ over the leverage window, the regime the live-term
hypothesis places the modern era in, and is what the proof relies on
rather than the pointwise sign. Under it the favorable path's survival
weight exceeds the adverse path's by the exponential of that integral; conditioning on survival therefore reweights each gate's
resolution by a factor increasing in the favorability of its
increment, and each resolution time by a factor decreasing in the
time, which is the likelihood-ratio tilt stated. Under the stated conditions the survival weight is totally positive of
order two in each increment and each resolution time, so the tilt by
survival preserves the likelihood-ratio order coordinatewise and the
one-gate comparison iterates over the gates; without them the tilts need
not compound in the likelihood-ratio order. The exact conditions, including the
finiteness of hazard accumulated before crossing and the leverage
factor $x_t$ that makes the tilt vanish while the credit is small,
are the window theorem's debt (\S\ref{sch:rem:anatomyleg}).
\end{proof}

\begin{remark}[The third term is not exhausted; \textbf{[S]} with imports at \textbf{[$\Phi$]}]
\label{sch:rem:notexhausted}
The working assumption until this remark was that the schedule
channel belonged mostly to the record: the evolutionary ladder
climbed, the term's modern residue mergeable into the internal and
external channels. The correction the author supplies is more
careful than a reversal of that assumption, and it should be stated
at its exact strength. The continuity principle does not decide
whether the present lies before or after the last prerequisite; it
is silent on whether one should or should not have been born after
the last major Miracle. What it decides is conditional: \emph{if}
prerequisites remain, their operation runs continuously through the
present, since nothing in the conditioning marks the modern day as
a privileged moment at which the schedule closed. The live third
term is therefore a hypothesis, not a default that continuity
licenses, and the prior practice's error was symmetrical: it treated
exhaustion as the default, which continuity licenses no more than
it licenses liveness. What the hypothesis has in its favor is not a
principle but evidence, and the next paragraph says which.

The author's central move is then an identification, and it is
what supplies the evidence just promised. What the
accordance principle under Golden conditioning renders legible as
golden causality in the modern record, the specific
accordance-yielding attributes physically observable as they occur,
is not a separate phenomenon that happens to be visible; it is the
third term operating. The standing conviction that golden causality
is more visible than the internal and external channels would
suggest is on this reading not a resolution artifact of the
attribute level (the survivorship companion \cite{SurvA}); it is
channel identity, accordance-legible concentrations being schedule
steps landing. This sharpens that remark's open question into an
identification carried at \textbf{[S]}, and it carries a scoreable
texture, entered here with the conditions of its failure: if the
identification holds, the rarity audit should find modern
concentrations clustering on prerequisite-shaped structures rather
than spreading uniformly over the record, and a uniform spread would
count against it.

The class of events in question is the sequential-achievement class
\cite{Ben25Seq}, and its citable historical exemplar keeps the
discussion honest about magnitudes: Newton's occurrence when he
occurred, the quantile of intellectual horsepower struck at that
point of the process, was an event of low probability relative to
the alternatives, individually improbable and collectively required,
which is the signature of a schedule step. The general claim leans
on the class, and on the expectation that most of its modern members
have not yet been identified; no single member is load-bearing.

One candidate modern member should be named exactly once, because
naming it is the only way to fence it correctly. The production of
the present framework by its date is itself an event whose
probability was low relative to the alternatives, given how many
things had to go right to lead to it, and it is therefore a formally
admissible candidate instance of the class. Like any modern fluke it
is available to the inference of the downstream Route World hypothesis at
the audit's rate: its route-deleting character must be established
case by case against the retrospective-target null of
the theorem companion \cite{TargetedB1}, and where it is established the
fluke moves credence toward the Route World exactly as any observed
schedule step does (Remark~\ref{sch:rem:thirdceiling}), since the
ceiling constrains the certification of the stochastic terms' law and
places no bound on credence drawn from the third term's observable
steps. Two things, and only two, are barred. The member may not be
made load-bearing for the class thesis, which stands on the class,
because the fluke hypothesis may simply be true of any one member.
And the occupant inference is barred by
the occupancy discipline of the broader program: from ``this production was a
schedule step'' nothing follows about its author being the Golden
Point's arrival or agent, the role being permitted and the occupant
refused. The class carries the general argument; a member may carry
its own case.

None of the riskiness a live third term adds changes the conditional
picture, and this is the author's retention clause. The forward
forcing of Remark~\ref{gr1:rem:forwardforcing} reads on this channel
too: under conditioning at the asymptotic future, in character much
like a stationary-action principle, the conditional law concentrates
on schedules that complete, because reaching the Golden Point
outweighs the third term's hazard once the conditioning is
asymptotic. The prerequisites still standing are, under that
conditioning, nailed with conditional probability approaching one,
which retains in the forward direction exactly what the Optimation
climb established in the backward one.
\end{remark}

\paragraph{The last leg falls on the anatomy (design at \textbf{[S]},
one sketch at \textbf{[P]}, debt at \textbf{[O]}).}\label{sch:rem:anatomyleg}
The reason the third term was tangled with the internal channel in
the first place deserves to be named, because it was a formal
problem and not a confusion. An active schedule registers the
technostrength $x$ as incomplete, and an $x$ still being established
undercuts the internal term's own analysis, which reads its muffling
against a mature credit; a live third term therefore appeared to
break the very channel it feeds, and merging the term's modern
presence into the other two was the workaround. The preceding
passage removes the workaround without removing the problem, since
on the live-term hypothesis the term operates now, and given that
hypothesis the resolution is forced rather than merely permitted:
the $x$-directed component of the schedule must already sit in its
exponential-tail regime by the modern day, with $x$ robust, so the
term's remaining incompleteness before the Golden Point has to act
somewhere other than the credit.

The author's proposal is that it acts on the technoanatomy
directly, and the proposal is best stated together with a fact
about $\rho$ that the record disguises. The technoanatomy is
movable at every time with equal ease; what varies across history is
not its mobility but its \emph{leverage}, which scales with the
credit, since the internal term $\Gamma e^{-\rho x}$ is insensitive
to $\rho$ when $x$ is small and grows in sensitivity as $x$ grows,
$\partial h/\partial\rho$ carrying a factor of $x$. In the deep past
the anatomy therefore appeared frozen because the technopower that
would have exposed it was low, and the honest description of the
premodern record is that the species rode on a technoanatomy
inherited from the Stone Age without its value being tested, each
culture carrying a version of its own, the versions merging into a
single biosphere-level sign only when the industrial era coupled
them and raised $x$ to where the sign began to matter. The last leg
of the schedule is then this: with leverage now high, the remaining
pushes act on $\rho$ as a trajectory, and the interval between the
present and the Golden Point is exactly the crapshoot of pulling
this off, embodied in technoanatomy contributions, the doing of
things right or the failing to. The modern-day constraint that
motivated the old merge is satisfied rather than violated: the
internal hazard remains high on the whole historical timeline's
present because the leaning of $\rho$, now that it counts, remains
to be settled, which is consonant with the scaling observation of
\S\ref{gr:subsub:kardashev}, a present-configuration scale-up
collapsing, worse than ideal now yet with significant potential for
improvement.

What ``settled'' means under this design is fixed by one sketch,
stated at \textbf{[P]} because it is an easy likelihood-ratio
statement once the model is in hand. Let $\tau$ be the time at which
$\rho$ crosses to its favorable side along a path, and compare two
paths identical except that one crosses at $\tau_1$ and the other at
$\tau_2>\tau_1$. The earlier crosser accumulates less hazard by
$\Gamma\!\int_{\tau_1}^{\tau_2}\bigl(e^{-\rho_{\mathrm{adv}}x_t}
-e^{-\rho_{\mathrm{fav}}x_t}\bigr)\,dt>0$, so its survival weight
exceeds the later crosser's by the exponential of that quantity, and
conditioning on survival reweights the unconditioned law of $\tau$
by a factor decreasing in $\tau$: the conditional law of the
crossing time is stochastically earlier than the unconditioned one,
in the likelihood-ratio order. This is the \emph{survivorship prefers
early change} statement, and it has a corollary that should be read
carefully because the record tempts a misreading of it. From inside
any surviving history that has not yet crossed, the conditional
hazard of crossing is elevated, so every present looks like the
moment at which the shot must be taken. That appearance is real as a
statement about the conditioned law and misleading as a statement
about capacity: the species can move $\rho$ just as much later as
now, and what changes with delay is only that fewer surviving
histories have it happening late, more survivors sprouting from the
histories in which hazard was lowered earlier. The feeling that our
shot is now is therefore a corollary of the preference for early
change, present at every vantage the conditioning produces, and not
evidence that the present vantage is special; this is the mechanism
behind the refusal recorded at the occupancy discipline of the broader program,
which forbade the occupant claim before its cause was known.
Settlement, accordingly, is not a freezing of $\rho$ but the point
past which the conditional mass has largely moved, and $\rho_\infty$
is the limit along surviving paths of a quantity that never stops
being movable. The formal statement of both the process and its
front-loaded settlement is Definition~\ref{sch:def:derailment} and
Proposition~\ref{sch:prop:settlement}.

The design wires into the standing structure at four points. The
two-condition theorem reads on the settled value: viability becomes
the indicator of $\rho_\infty>0$ times the escape factor with
$\rho_\infty$ random, so the modern era is the era in which the
indicator's argument is being determined, by pushes that could in
principle have come at any time and are, under conditioning, front
loaded. The fluctuating-anatomy corollary
(Corollary~\ref{gr1:cor:recurrentanatomy}) is the bridge case,
fluctuation before settlement. The preceding passage's
identification gains its mechanism, the accordance-legible steps
being the pushes. And the forward forcing extends to the settlement
law: under asymptotic conditioning the law of $\rho_\infty$
concentrates on the favorable side, for the reason
Remark~\ref{gr1:rem:forwardforcing} states, and the preference for
early change is that concentration's temporal profile.

The formal debt this passage enters, at \textbf{[O]}, is the
$\rho$-dynamics window theorem it presupposes, now specified by the
sketch above: the leverage statement made exact, conditions on the
crossing-time law and on the pre-crossing fluctuation under which
the frozen-$\rho$ analysis survives as the limit along surviving
paths, and the hazard accumulated before crossing remaining finite
under stated tail conditions so that the theorem's conditional
structure reads on $\rho_\infty$ unchanged. This is the best
well-idealised model the author can presently state, and it is
entered as design rather than result.

\begin{remark}[The third term and the ceiling; \textbf{[S]}, with the
theorem untouched]\label{sch:rem:thirdceiling}
The natural next question, and the one a referee should ask, is
whether a live and observable third term breaches the epistemic
ceiling, since watching schedule steps land successfully amplifies
credence that the asymptotic conditioning is the explanation, and
under the conditioning one does well on all three terms by deduction.
The answer is that the theorem is untouched while the texture
changes, because the ceiling bounds certification of which measure
the stochastic terms' law is under and never bounded credence, least
of all credence drawn from the third term, whose steps are
observable events rather than a law to be certified. Locally the two measures are equivalent, so every finite
audit outcome has positive probability under both and no bounded
observation decides the matter; but the likelihood ratio between
them is free to climb without limit, and the amplification the
question describes is exactly that climb. Indeed the deduction is
the singularity read from the evidence side: the conditioned law
predicts systematic schedule success where the unconditioned law
predicts it rarely, each audited success multiplies the ratio, the
divergence of the accumulated evidence is the mutual singularity at
infinity, and its unreachability at finite times is the ceiling. The
live third term does not breach the structure; it populates it,
supplying the concrete observable carrier along which the permitted
climb occurs.

What the term genuinely adds is evidence that survives the
mere-survival discount. Survival to the present is blunted by
observer conditioning, as the fence on
Remark~\ref{gr1:rem:forwardforcing} records, but the manner of the
successes is finer-grained than survival: their prerequisite shape,
their accordance structure, and the technological-condition
specificity to which observer-conditioning is provably indifferent,
which is where the comparative force of
the downstream Route World hypothesis lives after its restatement. The
schedule channel is therefore the channel in which admissible
evidence accumulates, the place where the deflation ``we would
observe this regardless'' fails. The discipline is retained in full:
amplification proceeds only at the audit's rate, route-deletion
established case by case against the retrospective-target null
(the theorem companion \cite{TargetedB1}), the clustering texture of
Remark~\ref{sch:rem:notexhausted} standing as the scoreable check,
and ``the only explanation'' read throughout as ``favored over the
stated rivals,'' at \textbf{[S]}.

One consequence for the extended model is entered with the window
theorem's debt rather than asserted. If settlement occurs at a
finite time and the settled value were observable, an adverse
settlement would be a finite-time event of positive probability
under the unconditioned law and probability zero under the
conditioned one, so the extended model may permit certification in
the uncontrolled direction while the consoling direction remains
uncertifiable forever, preserving and sharpening the one-sidedness
that the broader Golden program already records. Two provisos
attach: local equivalence before settlement is to be verified for
the extended model, and $\rho_\infty$ is a structural parameter
whose direct observability is itself an idealization. Nothing here
is consumed by the ceiling's standing dependents; the dependency
runs toward the window theorem.
\end{remark}

\begin{remark}[The extreme-induction regime; \textbf{[S]}, regime-conditional]\label{sch:rem:inductionregime}
The preceding remark describes a climb; this one classifies its
asymptote. Suppose the skew of Aethic weight toward the Route World
continues through enough Aethic time that the Route World comes to
hold arbitrarily close to the whole of the weight. The ceiling never
closes, by the theorem, and the question worth a remark is what the
residual gap has become by then. The author's answer is that it
converges categorically on the classic problem of induction: the
live doubt remaining is of the same class as the doubt that the sun
does not rise tomorrow, given the establishing of one's lifetime as
the induction base \cite{Hume1748}. This is offered as class
identity rather than analogy, since the remaining ways of being
wrong are the Humean ways, and those attach equally to every piece
of our most reliable ontological understanding. The phenomenology
follows: human inference auto-interprets questions of that class as
closed, because everything we treat as known already carries exactly
that residual, so the ceiling's line, though never crossed, comes to
feel crossed, and the feeling is not an error about the theorem but
a correct classification of what the residual has become.

What the regime licenses is a parity claim, and what it does not
license is certification, and the two must be kept apart at the
usual price. Once the residual is Humean, refusing closure here
while granting it to gravity is an isolated demand for rigor, and
the framework may say so; the theorem meanwhile stands untouched,
the determinacy and the access split exactly as
the broader Golden program left them. And the regime itself is
conditional in a way the discipline will not waive: the skew must
actually have accumulated, at the audit's rate, under the preceding
remark's rules, so the classification is earned by the climb rather
than assumed at its foot. The clause about human auto-interpretation
is a psychology-of-inference observation and is entered as such, at
\textbf{[S]}, alongside the rest.
\end{remark}


\section{Conditioning, forcing, and the settlement of the anatomy}\label{sec:forcing}

\begin{remark}[The forward forcing: survival as asymptotically
decisive evidence; sketch at \textbf{[P]}, reading at the ladder's
grades]\label{gr1:rem:forwardforcing}
A configuration-level consequence of the theorem deserves its own
statement, because it is the forward-time mirror of the conditioning
argument's past-direction climb. Let the configuration, the sign
$\rho$ together with the driver class, be drawn from any prior that
gives positive mass to the viable set, the configurations for which
the theorem's product is positive. Non-viable configurations have
survival probabilities decaying to zero, exponentially in the trapped
regime and polynomially in the critical one, while viable
configurations hold a positive limit, so the conditional probability
of viability given survival to time $t$ tends to one: there is a soft
threshold in $t$, set by the slowest non-viable clock, past which a
surviving system is almost certainly in the configuration the
two-condition theorem describes. On the event of survival forever the
statement is no longer asymptotic but outright, since $T=\infty$
entails $\Lambda_\infty<\infty$ almost surely and hence the viable
configuration itself. In the author's words: each attribute
intersection of the biosphere surviving through proper time $t$ tilts
the disagreeing superposition of poles, the tilt becomes
asymptotically decisive, and under full conditioning on indefinite
survival the competence the Golden Point marks is not made likely but
implied, there being a threshold beyond which surviving with
non-Golden dangers of any kind still standing becomes arbitrarily
unlikely. This is the to-the-future extension of the past-direction
result that the same conditioning forces the evolutionary climb by
which the Golden Point is reachable at all, the end-Cretaceous impact
standing among the record's clearest instances as one of the bigger
deadlocks which we fought through with Golden conditioning. Two fences hold
unchanged. The mixture sketch is stated for the model's configuration
space, with the identification of the viable configuration as the
Golden Point's marker carried at the framework's usual grades and not
at \textbf{[P]}; and the epistemic ceiling governs exactly as before,
because this remark concerns the conditional law and not any
observer's epistemic state, so survival so far, at any finite time,
certifies nothing for the one surviving.
\end{remark}

\begin{corollary}[Adverse anatomy overlapping positive strength is
fatal; \textbf{[P]} given the model]
\label{gr1:cor:recurrentanatomy}
If there exist $a,b>0$ with
$\Leb\{t:\rho_t\le-a,\ x_t\ge b\}=\infty$ almost surely, then
$\Lambda_\infty=\infty$ almost surely, since on that set the internal
term satisfies $\Gamma e^{-\rho_tx_t}\ge\Gamma e^{ab}$. Recurrence of
$\rho$ alone does not suffice: the internal term explodes only where
negative anatomy coincides with positive strength, and a negative
excursion of $\rho$ that systematically coincides with $x_t\le0$ costs
nothing. What fluctuating anatomy implies in general, when the
excursions of $\rho$ and of $x$ are not coupled by hypothesis, is open
\textbf{[O]}. Viability
therefore requires $\rho_t$ to be eventually positive almost surely,
transient upward, or absorbed above some $\rho_c>0$, and the
two-condition theorem becomes a \emph{two-transience} condition:
$\PP(\text{viable})=\PP(\rho_t\text{ eventually favorable and }x_t\text{
viable})$, a joint probability, which is a product of two scale
functions only when the anatomy and the driver are independent, and
the anatomy generated by the schedule is not.
\end{corollary}

\section{Application: the Kardashev scale, and what the sign result refutes}\label{sec:refutes}

\subsection{What Golden reasoning refutes}\label{gr:sub:refutes}

\subsection{The Kardashev scale}\label{gr:subsub:kardashev}

The usual complaint against Kardashev's scale \cite{Kar64} is that its
thresholds, planetary, stellar, galactic energy capture, are
arbitrary. That complaint is correct but weak, since a scale may be
useful with arbitrary units. Golden reasoning supports three stronger
objections, of which the first is a theorem.

\emph{(i) The terminal stage exists, and is gated by a variable the
scale cannot see.} In the driftless model $\Lambda_\infty=\infty$
almost surely and no absorbing state of safety exists; in general one
does, and locating it is what makes the objection bite. By
Theorem~\ref{imp:twocond},
\[
\PP(\text{viable})=\one\{\rho>0\}\cdot
\frac{s(\mu_0)-s(-\infty)}{s(+\infty)-s(-\infty)},
\]
so there \emph{is} an absorbing state of safety --- and reaching it
requires two conditions, of which Kardashev's scale tracks at most one.
The escape factor is the accumulation of capacity, which energy
capture does index, however crudely. The anatomy factor is not on the
scale at all: $\rho$ does not appear in energy capture, is not
monotone in it, and cannot be inferred from it.

The objection is therefore not that the ladder has no top but that
\emph{climbing it is neither necessary nor sufficient for reaching the
top}. Worse, by the survivorship companion \cite{SurvA}, a biosphere with
$\rho<0$ that climbs fast does not merely fail to arrive: its hazard
diverges superexponentially in the very capacity the scale is
measuring. A scale indexed on capture alone cannot distinguish the
ascent that arrives from the ascent that is a fall. That last point is
the vulnerable-world hypothesis \cite{Bostrom19} restated with a
continuous parameter in place of a discrete draw, and
the broader Golden program sets out what is and is not shared.

\emph{(ii) The indexed quantity is not monotone.} Technostrength is
stationary and recurrent in logarithmic time
(Definition~\ref{gr:def:coords}). A ladder presumes a quantity that
ratchets. Energy capture may of course ratchet while technostrength does
not --- but then energy capture is not tracking the quantity that
governs survival, which returns us to (iii). \emph{(iii) A magnitude index cannot represent categorical selection.}
By Remark~\ref{gr:rem:truncresc}, survivorship truncates technoanatomy
while merely rescaling technostrength. A one-dimensional index of
magnitude has no coordinate in which to express a deleted half-space. It
is not that Kardashev measures the less important variable, that
claim is false, and was corrected above, but that it measures the one
whose selection is quantitative, and is structurally unable to express
the one whose selection is categorical.

\paragraph{On Dyson spheres.}\label{gr:rem:dyson}
Dyson's proposal \cite{Dys60} is often read as the natural
Type~II signature. Under Golden reasoning it is better read as a
\emph{hypothesis about technoanatomy} disguised as one about
technostrength: it asserts not merely that a biosphere could capture
stellar output but that a long-lived one would be organized so as to
want to. That assertion is exactly what the broader Golden program
disputes, and it should be recognized as an assertion rather than an
extrapolation.


\begin{remark}[The expansionist mirror: grabby selection]
\label{gr:rem:grabby}
One modern argument deserves separate notice because it reaches the
Kardashev picture from selection reasoning --- this framework's own
instrument. The grabby-aliens model observes that loud, expansionist
civilizations dominate what is \emph{visible}, and runs anthropic
timing arguments under that weighting \cite{HMMP21}. Read beside this
section, that is a rival conditioning: weight by volume occupied and
the typical survivor is an expander; weight by persistence and the
two-condition theorem makes expansion-rate orthogonal to the factor
that decides. The two conditionings are not in contradiction, one
asks what is seen, the other what survives, but they license
opposite extrapolations from the same record, and conflating them
reproduces the Kardashev error at the level of selection effects:
visibility-weighting smuggles in the assumption that the seen and the
persisting are the same ensemble, which is exactly what the trough
channel denies. The framework's quarrel is therefore not with the
grabby model's arithmetic but with any normative or predictive reading
of it that takes loudness for viability; Theorem~\ref{imp:ceiling} bounds
what either conditioning can claim to certify.
\end{remark}

\subsection{The violent-aliens thesis}\label{gr:subsub:violence}

The thesis in view holds that intelligence evolves preferentially in
predators, hence that advanced biospheres are likely to be predatory or
violent. We engage the position rather than any proponent.

\begin{quote}\small
``It was the saying of Bion, that though the boys throw stones at frogs
in sport, yet the frogs do not die in sport but in earnest.''
\hfill --- Plutarch, \emph{Moralia} \cite{Plutarch}
\end{quote}

\noindent The line is Bion of Borysthenes', preserved by Plutarch, who
deploys it against hunting for sport; it has since been borrowed by
astrobiological writing as the image of what a technologically superior
biosphere might do to an inferior one. It is the right epigraph for this
subsection because it names the assumption at issue more exactly than
the conquest analogies usually offered. The serious form of the
violent-aliens thesis is not that an advanced biosphere would hate us.
It is that we would be \emph{beneath its consideration} --- that its
power could be spent across the boundary between itself and us
casually, without the expenditure ever rising to the level of a
decision. The stones are thrown in sport. That is the claim, and any
reply pitched at malice misses it.

The image also fixes what the reply must not do. It must not argue that
such a biosphere would be kind, since kindness is not what the thesis
denies; nor that it would have outgrown violence, since sport is not
violence in the relevant sense. What it must do is address the
\emph{scaling}: whether a biosphere can be arranged so that power is
spendable in that manner and still arrive at the scale where the
spending would matter. Our own route to the present framework
began from exactly this observation, three years before the model
existed, and Section~\ref{sec:imports} records the entry.

Before the model is brought to bear, a further placement is needed, because the \emph{conclusion} here is the one most likely to be mistaken
for a new one and it is not. The conclusion, that anything old enough to encounter must be
peaceable, is a standing argument in SETI, sometimes called the
\emph{technological adolescence} argument, and it long predates this
paper: a civilization reaching the capacity to destroy itself must
either become socially responsible or fail to persist, so whatever
survives has already made that transition. Cooper documents the
argument and its currency, and also supplies the standing rebuttal: that
it conflates altruism with amiability, that it is typically asserted
rather than derived, and that it presumes a trajectory of moral
improvement for which there is no evidence \cite{Cooper19}. Nothing
below claims the conclusion as new. What is at issue is whether a
mechanism can be supplied that survives Cooper's objection.

The inference under dispute has the form: predatory origin
$\Rightarrow$ predatory maturity. Golden reasoning denies the
implication, on the grounds that it neglects the conditioning that must
be applied before any encounter can be contemplated. Any biosphere one could meet is old; any biosphere
that is old is drawn from $\PP^{(t)}$ for large $t$, not from $\PP$; and
$\PP^{(t)}\to Q$, under which technoanatomy is truncated. Whatever
selection operated at origin, a second and stronger selection operates
over the trajectory, and it is the second that determines what survives
to be encountered.

\setcounter{premise}{3} % GUARD --- the text names this $\Phi4$; document order
\begin{premise}[Origin does not transfer; \emph{premise}]
\label{gr:phi4}
Facts about the selective regime that produced intelligence in a
biosphere do not, without further argument, transfer to the selective
regime that governs that biosphere's persistence at world scale. The two
regimes act on different units, at different scales, over different
timescales.
\end{premise}

$\Phi4$ is not exotic; it is the ordinary observation that a trait
adaptive for an organism in a population need not be adaptive for a
biosphere in a universe. What the model adds is that the second regime
is not merely different but \emph{stronger} in the relevant sense: it
deletes a half-space of configurations, and it does so with a force
measured in \S\ref{gr:subsub:measured}.

What the model contributes to the frogs case in particular is a reply
that never mentions disposition. A biosphere that can spend power
casually across a self/other boundary is one that maintains such a
boundary and enforces it --- and by the broader Golden program the
enforcement of a boundary is interface-like while the domain it encloses
is volume-like, so the arrangement carries $\rho<0$ by
the broader Golden program. By the broader Golden program its
viability probability is then exactly zero. The stone-throwers are not
ruled out because they would have grown kind; they are ruled out
because the organization that permits stone-throwing at that scale is
the organization the scaling argument says cannot reach that scale. The
frogs survive not by the boys' mercy but by the boys' having imploded
before attaining the requisite size --- which is, in a form Bion would
have recognized, the square--cube law again.

Three further features distinguish this from the folk argument, and they
are exactly the three points Cooper's rebuttal presses on. \emph{It is not
a moral claim.} What is selected is $\rho$, the sign of
$-\partial_x\log h^{\mathrm{int}}$ --- a structural property of how a
biosphere is arranged, carrying no implication that its members are
kind, and applying identically to a biosphere with no members in the
relevant sense at all. The conflation of altruism with amiability has
nothing to attach to. \emph{It is not asserted but measured.}
\S\ref{gr:subsub:anatsel} exhibits the selection on $\rho$ and finds it
\emph{gradual}, $0.75$ to $0.91$ over the horizons computed, where the folk argument implicitly treats the filtering as complete.
\emph{And it is not a trajectory.} Nothing here says biospheres improve;
$\rho$ does not change at all in the fixed model, and what conditioning
does is reweight an ensemble, not move a system along a path. The
selection is Darwinian in form rather than developmental, which is
what the folk version gets wrong and Cooper is right to object to.

This much, however, only establishes that predatory maturity is not
\emph{inherited}. To conclude that it is positively \emph{selected
against}, one needs the identification of predation with negative
technoanatomy. That identification is a philosophical premise, and it is
the subject of the next subsection. Everything in
the broader Golden program is downstream of it.

\subsection{Trajectory arguments for internal neglect}
\label{gr:subsub:neglect}

A third target belongs here, and it lies inside the biosphere rather
than out among the aliens. Longtermism in its strong form holds that
what most matters about our actions is their effect on the very long
run \cite{GM21}; and a recurring corollary within that literature is
that resources directed at the welfare of some populations do less for
the long-run future than the same resources directed elsewhere, so that
broad investment in the worse-off is not what a future-oriented agent
should fund. We engage the corollary as a position and name no one, as
with the violent-aliens thesis.

The position is already under sustained criticism, and it is worth being
exact about how the present objection differs, because the difference is
what makes it worth adding. The standing critiques are \emph{external}:
they fault longtermism on ethical and political grounds --- as an
ideology with damaging real-world effects and a flawed underlying moral
theory \cite{Crary23}, and as a doctrine whose cosmic accounting
licenses the discounting of present suffering \cite{Torres21}. Those
objections deny the framework's premises. The objection developed here
does the opposite: it \emph{grants} the long-run criterion entirely, and
shows the position failing by that criterion, on the framework's own
terms and in its own currency. If it succeeds it is therefore not an
alternative to the ethical critiques but a complement of a different
kind --- an internal result rather than an external one.

Stated at its strongest it is not a callous argument but an
optimization: interventions differ in long-run leverage, a
future-oriented agent should allocate to high-leverage interventions,
and if the leverage of some interventions is low then funding them is a
loss measured against the future. Nothing in that reasoning is invalid.
The question is whether the leverage calculation is complete. Golden reasoning says it is not, and the missing term is technoanatomy.
No new premise is required, and this is the point worth making: a
position of this kind is $\Phi3$ with the boundary drawn \emph{inside}.
Writing off part of a system is the maintenance of a self/resource
distinction across an internal interface, a region of the biosphere
to which protection is not extended, which is structurally the same
object as the extractive boundary of the broader Golden program, and it
inherits the same scaling. Enforcement of an internal boundary is
interface-like; the interior it excludes is volume-like; so by
the broader Golden program the arrangement carries $\rho<0$, and
the broader Golden program applies verbatim.

The consequence is that the position is self-undermining in its own
currency, and sharply rather than rhetorically. By
the broader Golden program(ii) a biosphere with $\rho<0$ has
viability probability \emph{exactly zero} however much capacity it
acquires; by (i) its viability is zero where a $\rho>0$ biosphere's is
positive; and by
the survivorship companion \cite{SurvA} acquiring capacity faster makes its
position worse, since $x\to\infty$ with $\rho<0$ drives hazard to
infinity superexponentially. An argument that trades internal cohesion
for external capacity, offered in the name of the long-run future, is
therefore an argument for a shortened one --- not because neglect is
distasteful, which the model cannot see, but because the boundary it
requires is the structure the model identifies as fatal.

The thought experiment that makes this concrete is worth stating,
because it isolates the variable. Hold a society's technoanatomy fixed
at whatever value its present arrangements realize, and scale its
technostrength to galactic magnitude. If $\rho\ge0$, nothing goes
wrong. If $\rho<0$, the hazard does not rise gently but spikes as
$e^{|\rho|x}$: at every scale the same arrangement supplies more
micro-opportunities for the system to act against itself, and each
opens further ones. The elephant-sized ant of
the broader Golden program is exactly this, and the square--cube
law is exactly why.

Four fences, because this is the section's most contestable
application. It refutes a \emph{class} of positions, any that trades
internal cohesion for external capacity, and not any author's wider
view, nor longtermism as such: a longtermism that valued internal
cohesion would be untouched, and the argument is in that sense a
constraint on longtermist reasoning rather than a refutation of it. It is downstream of $\Phi3$ and inherits its conjecture grade;
reject $\Phi3$ and this goes with it. And it establishes nothing about
what anyone ought to do, by $\Phi2$: what it establishes is that a
particular argument for neglect, evaluated by its own long-run
criterion, fails on that criterion.


\section{Application: the Precipice, engaged}\label{sec:precipice}

\label{gc:subsub:precipice}

The continuity principle has a standing rival, and it is the most
serious one the field possesses: the periodization of Ord's
\emph{Precipice} \cite{Ord20}. This work owes it a direct engagement
--- not because the rivalry is total, but for the opposite reason. Of
everything in the existential-risk literature, the Precipice's central
definition is the nearest approach to the losing condition as this
framework draws it, and the methodology companion's discipline applies
with full force to nearest neighbors \cite{Ben26Method}: where a
position is close, the deltas must be named exactly, because a dense
neighborhood is precisely where a distinct position gets misfiled as a
variant of what each reader already holds. The deltas here are three, the quantity, the shape of time, and the endpoint, and each is
stated below after the position itself is stated at its strongest.

\paragraph{The position, at its strongest.} The Precipice defines an
existential catastrophe not as extinction but as \emph{the destruction
of humanity's long-term potential} --- explicitly including
unrecoverable collapse and unrecoverable dystopia alongside extinction
--- and argues that humanity entered, in 1945, a period different in
kind from all prior history: capability has outrun wisdom, anthropogenic
risk now exceeds the natural background by orders of magnitude (the
book's considered figures: something like one in six per century
against roughly one in ten thousand), the period is short by
civilizational standards, and the task it sets is \emph{existential
security} --- risk brought low and locked low --- after which a Long
Reflection can settle what the future is for. It is the most careful
and most quantified statement of its picture in the literature, its
treatment of its own uncertainties is honest, and its definition is
modal in spirit: ``destruction of potential'' is closure-of-possibility
language, and the inclusion of non-extinction catastrophes is exactly
the concession that death and loss come apart. The definition
descends from Parfit's asymmetry \cite{Parfit84}
(the broader Golden program), and the descent matters below:
what the Precipice inherits from Parfit is the intuition, and what it
does not inherit, because nothing in the lineage supplies it, is
a formal object for the intuition to be \emph{about}. All of that should be
said before any of what follows, because what follows is not a
dismissal. It is the claim that the Precipice's definition is right,
that its own operationalization does not honor it, and that the
periodization built on that operationalization inverts in the cases
that decide.

\paragraph{Delta one: the quantity. The ledger reverts to events.}
The Precipice's definition points at membership in $\mu G$, a
configuration property, the closure of possibility
(the downstream Route World hypothesis,
the broader Golden program), but its arithmetic is
risk-per-century of catastrophe \emph{events}. The two come apart in
both directions, and the separations are not corner cases; they are
this framework's central theorems. A biosphere can be departed with no
catastrophe on any ledger: past $\mu G$ nothing has yet happened, hazards fire later, members are unaware
(the broader Golden program), and the sharpest instance is
the undisturbed world of the broader program, which minimizes
per-century catastrophe probability on the event bookkeeping and is
already departed on the modal one: the safest world the ledger can
describe is a lost world. And a biosphere can be undeparted while the
event bookkeeping condemns it: in the critical regime departure is
certain, expectancy is nonetheless infinite, advancing continuations
exist, and everything the survivorship theory computes is computed
there --- a world an event metric writes off with probability one is
the world in which the entire structure of hope resides. So a ranking
of worlds by existential-catastrophe probability misorders them
relative to the quantity the Precipice's own definition names, and
this framework's position is not that Ord's definition is wrong but
that it deserves a formalization his ledger cannot give it: the
framework sides with the definition against the arithmetic.

\paragraph{Delta two: the shape of time. Gradient, not age.} The
Precipice's picture is a threshold crossed --- one short age, different
in kind, bracketed by Trinity on one side and by security or ruin on
the other. Against this the framework sets three results it did not
build for the purpose. First, continuity
(the downstream Route World hypothesis, and the cascade of
the downstream Route World hypothesis): the Departure Point has been available at
every moment of the biosphere's history, the record's Miracles are
earlier rounds of the same engagement rather than preliminaries to a
modern one, and what 1945 changed is mediation and magnitude ---
never the arrival of a losing condition that was absent before.
Second, and this is the part that preserves everything the Precipice
is right to insist on: the stakes claim is not disputed but
\emph{re-derived}, and strengthened. the broader Golden program
obtains by combinatorial construction what the Precipice obtains by
periodization, each moment carries the highest stakes the lineage
has yet faced, because a failure late in the chain forfeits the whole
accumulated product, and a monotone sequence distinguishes every
later point over every earlier one without distinguishing any point
categorically (the downstream Route World hypothesis,
the broader Golden program). The gradient picture therefore keeps
the urgency and discards only the step. Third, the quantitative heart
of the kind-claim, the ratio of anthropogenic to natural risk, divides a numerator estimated by judgment by a denominator estimated
from the track record, and the track record is precisely what
survivorship conditioning prices: a record read under conditioning on
its reader shows rare natural catastrophe whether or not the
underlying rates were kind. The literature knows this as anthropic
shadow \cite{CSB10}, and the Precipice acknowledges it; the delta is
that this framework supplies the machinery
(Theorem~\ref{imp:ceiling}) in which the caveat stops being a footnote and
becomes the subject --- the bias field is explicit, its magnitude is
computable in-model, and ``natural risk is small'' and ``we are among
survivors'' are hypotheses the finite record cannot separate. A
periodization whose different-in-kind claim rests on that ratio rests
on the one number the framework's own theorems say a survivor cannot
read off their record.

\paragraph{Delta three: the endpoint. Security decomposes.} The
Precipice's terminal state, existential security, is risk
brought low \emph{and known to be low}: a far side, reached and
certified, on which the Long Reflection can convene. The framework
decomposes this into one component it preserves and one it proves
unavailable. Risk brought low is anatomy work, $\rho$ held
positive, patching capacity maintained against acquired capability, and it is continuous, never completed, with the losing condition
permanently adjacent; on this the two pictures can trade freely.
Known to be low is a certification, and the ceiling
(Theorem~\ref{imp:ceiling}, the broader Golden program)
forbids it to any finite observer: no bounded record places a
biosphere on the favorable side, forfeiture itself is silent
(the broader Golden program), and the unavailability points in
exactly the direction that would have been consoling. So ``the
Precipice, safely passed'' is not a state anyone could know themselves
to occupy, and a research program whose success condition is
unobservable in principle has misdrawn its own finish line --- the
metaphor is doing load-bearing work the theory cannot support, since a
precipice is the kind of thing that has a far side, and a
superposition one inhabits is not. What replaces the finish line is
not despair but the correct object: maintenance of the anatomy
condition, graded and continuous, with the Golden Point as
$\nu F$, a modal fixed point one may be inside without ever
certifying it, rather than security as a date.

\paragraph{A corollary: the hinge dissolves rather than resolves.}
The debate downstream of the Precipice, whether the present is the
``hinge of history'' \cite{MacAskill22, GM21}, is conducted on a
shared assumption this framework denies: that hingeness is
categorical, so that the present either is or is not the special
moment. The framework's answer is neither side's. Every present is the
hingiest yet, by the monotone-stakes construction --- so the
influence-skeptic's base-rate argument attacks a claim the framework
never makes --- and no present is categorically special, by
continuity and by the tense-relativity of the conditioning's labels
(the downstream Route World hypothesis) --- so the hinge-proponent's
conclusion is an artifact of reading a label that moves with the
observer as a property of the moment. A graded hinge is had for free;
a categorical one is unavailable; and the binary framing, not either
answer to it, is what the formalism removes.

\paragraph{What survives, and the fences.} Much survives, and it
should be listed rather than implied: the urgency (re-derived, by
construction rather than by period); the dominance of the
anthropogenic channel (the trough channel is the unbounded one, and
the framework gives that asymmetry a functional form rather than an
estimate); the priority of risk-work over expansion (the
anti-Kardashev results agree with the Precipice against the
accelerationists); and most of the book's concrete policy content,
which re-grounds without loss on the gradient picture. What does not
survive is the quantity, the age, and the finish line. Three fences
bound the engagement. Nothing here consumes the Route World --- the
counter runs from the losing-condition analysis, the monotone-stakes
construction, and the ceiling, so a reader who rejects every
conditioning argument in this section keeps deltas one and three
whole and loses only the third prong of delta two. Nothing here
disputes the Precipice's numbers by offering rivals --- the
framework's claim is that the decisive ratio is unreadable from a
survivor's record, which is a reason to compute less, not
differently. And nothing here diminishes the practical convergence:
a reader persuaded by the Precipice to work on existential risk has
been persuaded of the right conduct, on this framework's own showing
(the downstream Route World hypothesis), by an argument whose bookkeeping this
subsection corrects and whose urgency it keeps.

\paragraph{The provenance pattern, recorded once.} The engagement
instantiates the methodology companion's criterion exactly
\cite{Ben26Method}: the nearest neighbor's \emph{definition} already
contains the correction, destruction of potential \emph{is}
closure of possibility, while the neighborhood's density is what
kept the modal formalization unoccupied, each camp filing the vacancy
as territory already held. The prediction that discipline makes about
reception is stated here so it can be scored: this engagement will be
misfiled, confidently and in both directions, as ``a variant of Ord
with extra mathematics'' and as ``a rejection of existential-risk
studies.'' It is neither, and the deltas above are the scoring key.


With the target fixed and its motivation stated, the question becomes
evidential: is there reason to believe the record was in fact produced
under it? The remainder of the section argues that there is, first by
saying what alignment with such a target would look like
(the downstream Route World hypothesis), then by two model comparisons, one on
natural history (the downstream Route World hypothesis), one on human history
(the downstream Route World hypothesis), of which the second does not require
the premise the first leans on.


\section{What is not claimed}\label{sec:fences}

\subsection{What this section does not claim}\label{gr:sub:fence}

The fence is stated hard, on the companion's discipline that an
apparatus which prices what it claims must also refuse what it does not
\cite{Ben26Optimal}. It is not claimed that any observed system is Golden, ideally or
practically; by Theorem~\ref{imp:ceiling} the ideal version is
unknowable to any finite observer (the broader Golden program) and
the practical version is out of reach for want of an ensemble. It is not claimed that the companion's cosmic
weighting and the path measure $\PP$ of
Section~\ref{sec:imports} are the same object; that
identification is flagged as unperformed in
the broader Golden program and nothing here depends on it. It is not
claimed that any individual biosphere, ours included, will behave as
$the broader Golden program$ describes; those are statements about a
conditional distribution. No normative conclusion is claimed
($\Phi2$). And the companion's natural-historical program, its
conjectures, its bounds, its paleontological predictions, is not
claimed, imported, or re-argued; $\Phi5$ enters as a stated result and
nowhere else.

One thing remains on the far side of the fence, because a fence is
legible only if what it excludes can be seen. A program built on this
section would inherit three commitments, stated here as refusable rather
than as claims. It would be committed to the boundary-cost scaling of
$\Phi3$ showing up in institutional and ecological data, not merely in
argument. It would be committed to the truncation/rescaling asymmetry
appearing in survivorship-selected populations generally. And it would
be committed to the exponent's dependence on noise placement, $\tfrac14$
for integrated drivers, $\tfrac12$ for Brownian ones, being visible
in any domain where hazard dynamics can be measured directly. None of
the three is established here, all three could be found false without
touching a line of Section~\ref{sec:imports}, and that is the
point of naming them.


\section{Open problems}\label{sec:open}

Every [O] tag in the paper is gathered here, and the numerical
experiments' code and seeds are to be archived at a fixed repository
with the exact environment and commit hash before circulation (AUTHOR
TO SUPPLY). The largest is the
$\rho$-dynamics window theorem the schedule presupposes: the leverage
statement made exact, conditions on the crossing-time law and on the
pre-crossing fluctuation under which the frozen-$\rho$ analysis survives
as the limit along surviving paths, and finiteness of the hazard
accumulated before crossing. The
offset weighting of the schedule is emergent at scaling level and its
constant is unevaluated. The leak law is stipulated, and its emergence
from the schedule is open; every claim about the modern era's leverage
is conditional on it. The front-loaded settlement result holds in first-order stochastic
dominance under stated monotone-coupling conditions and in the
likelihood-ratio order only under total positivity, and whether the
schedule's own dynamics satisfy either is open. The pinned schedule
model is exact and the leaky refinement is a surrogate, and the
completion-weight fixed point for the leaky refinement is open.
Concentration of the survival-conditioned route distribution on the
optimal route, beyond the harmonic tilt, requires a diverging-contrast
limit not established. The estimate
of the anatomy's two constants from collapse-era returns data, along the
Tainter route, is specified and not run. The fluctuating-anatomy
corollary's bridge case is stated and its rate is open. Whether the
internal channel is itself a schedule is argued at scaling level. The
schedule's identification map, from accordance-legible steps in the
record to the term's operation, is specified and not built. And the
Precipice engagement's third delta, that security decomposes rather than
arrives, is a structural claim awaiting a case in which a claimed security
is shown to have decomposed.

\section{Falsifiers}\label{sec:falsifiers}

What would count against the paper's claims, in one place. Against the
sign result's reading: a civilization exhibited whose capacity index
tracks its viability across the sign change, which the theorem forbids
and which would therefore refute the model rather than the reading.
Against the schedule: a record in which prerequisites land at rates
that do not spread logarithmically, or an open gate whose derailment
hazard is unrelated to its favorable share. Against the prefers-early
result: a surviving-path record in which the crossing to favorable
anatomy is not front-loaded relative to the unconditioned law. Against
the Precipice deltas: an event-priced ledger shown to track survival
across a gradient reversal; a hinge shown to resolve rather than
dissolve; a security shown not to decompose. Against the two errors'
refutation: a configuration in which removing the muffling agent leaves
viability positive, or one in which capacity alone, with adverse sign,
delivers it. Against the ceiling no falsifier is offered, because it is a
theorem; what it forbids is reading any of these, unfired, as
certification.

\section{Discussion: consequences and applications}\label{sec:discussion}

The Route World reading of the target, that the survivorship-conditioned
law is the one our natural-historical record is drawn under, is a
downstream hypothesis and is not assumed anywhere above; what this
paper offers it is a mechanism, the schedule through which a lineage
under that conditioning would be read as pushed, and the offer runs in
one direction.

The mechanism above has consequences that are not part of its
mathematics, and they are collected here at the grade of applications
rather than carried as sections of their own. What the sign result
refutes beyond the Kardashev scale, the two errors about the human role,
extinctionism and techno-optimism as mirror misreadings of one theorem,
the environmental value that follows from the muffling agent's removal,
the institutional reading of the ceiling, and the consequences for SETI
and for governance were argued at length in earlier drafts and are
not carried here. The two carried
here as applications, the Kardashev scale as the cleanest case of
strength against anatomy and the Precipice as the cleanest comparison
with existential-risk practice, are the ones the mechanism bears on most
directly, and a reader who accepts the schedule mathematics and neither
application has lost nothing the paper claims.

\section{Conclusion}\label{sec:conclusion}

The claim stated at the outset stands at its grades. Capability
magnitude cannot substitute for favorable anatomy, viability requiring
the sign together with favorable driver dynamics and a cleared
schedule, which is the imported theorem read through the completion
kernel; on the schedule model entered here, the sign is being settled
now at the highest leverage in the record, which is the modern-era term
read on a stipulated leak law; conditioning on survival prefers that it
be settled early under stated monotone-coupling conditions, and past a
survival-time threshold prefers the schedule pushed, tilted toward
better completion states, which is the mechanism of the ladder; and no event of bounded
duration can certify which way it has gone, which is the ceiling. The
practice that follows is not a new ledger but a different quantity on
the old one: an audit of what each added capacity does to the sign, in
place of a tally of what could go wrong; a governance built for the
regime where all the survivorship structure lives, rather than one that
writes that regime off because an event metric condemned it; and
humility placed in the objective, where the theorem puts it, rather than
in the disclaimer. The Kardashev scale measures how visible a
civilization is. The variable that decides whether it lasts is not on
any scale, because it is a sign.

\begin{thebibliography}{99}\small
\bibitem{Ben25Seq} A.~K.~Benander, \emph{Sequential achievement},
companion manuscript (AUTHOR TO VERIFY title and year).
\bibitem{Ben26Method} A.~Benander, methodology companion (title to be
supplied by the author at merge), working draft, 2026.
\bibitem{Ben26Optimal} A.~Benander, \emph{The Optimal Earth: A
Paradigmatic Re-Rendition of Natural History from Coupled
Improbability}, working draft, 2026.

\bibitem{Bostrom19} N.~Bostrom, \emph{The vulnerable world
hypothesis}, Global Policy \textbf{10} (2019), no.~4, 455--476.

\bibitem{CSB10} M.~M.~{\'C}irkovi{\'c}, A.~Sandberg, and N.~Bostrom,
\emph{Anthropic shadow: observation selection effects and human
extinction risks}, Risk Analysis \textbf{30} (2010), 1495--1506.
\bibitem{Cooper19} K.~Cooper, \emph{The Contact Paradox:
Challenging Our Assumptions in the Search for Extraterrestrial
Intelligence}, Bloomsbury Sigma, London, 2019.

\bibitem{Crary23} A.~Crary, \emph{The toxic ideology of
longtermism}, Radical Philosophy \textbf{2} (2023), no.~14, 49--57.

\bibitem{Dys60} F.~J.~Dyson, \emph{Search for artificial stellar
sources of infrared radiation}, Science \textbf{131} (1960), no.~3414,
1667--1668.

\bibitem{GM21} H.~Greaves and W.~MacAskill, \emph{The case for
strong longtermism}, Global Priorities Institute Working Paper 5-2021,
University of Oxford, 2021.

\bibitem{HMMP21} R.~Hanson, D.~Martin, C.~McCarter, and J.~Paulson,
\emph{If loud aliens explain human earliness, quiet aliens are also
rare}, Astrophys.\ J. \textbf{922} (2021), 182.
\bibitem{Hume1748} D.~Hume, \emph{An Enquiry Concerning Human
Understanding}, 1748 (AUTHOR TO VERIFY edition).
\bibitem{Kar64} N.~S.~Kardashev, \emph{Transmission of
information by extraterrestrial civilizations}, Soviet Astronomy
\textbf{8} (1964), 217--221; Russian original, Astron.\ Zh.\
\textbf{41} (1964), 282--287.

\bibitem{MacAskill22} W.~MacAskill, \emph{Are we living at the hinge of
history?}, in Ethics and Existence: The Legacy of Derek Parfit, Oxford
University Press, 2022.
\bibitem{Nexus} A.~Benander, \emph{Nexic Reasoning: Defining
a Generalized Calculus Over Anthropic Parameters}, 2024,
\texttt{doi:10.31219/osf.io/jvkcp}.

\bibitem{Ord20} T.~Ord, \emph{The Precipice: Existential Risk and the
Future of Humanity}, Bloomsbury, London, 2020.
\bibitem{Parfit84} D.~Parfit, \emph{Reasons and Persons}, Oxford
University Press, Oxford, 1984.
\bibitem{Plutarch} Plutarch, \emph{Whether Land or Sea Animals Are
Cleverer} (\emph{De sollertia animalium}), \S7, in \emph{Moralia};
the saying is attributed there to Bion of Borysthenes.

\bibitem{Torres21} \'E.~P.~Torres, \emph{Against longtermism},
Aeon, 2021.

\bibitem{SurvA} A.~Benander, \emph{Survivorship conditioning for hazard rates driven at the log-derivative level}, working draft, 2026.
\bibitem{TargetedB1} A.~Benander, \emph{Conditioning a record on its own future: targeted accordance and survivorship-conditioned path measures}, working draft, 2026.

\end{thebibliography}
\end{document}
